% generated by GAPDoc2LaTeX from XML source (Frank Luebeck)
\documentclass[a4paper,11pt]{report}

\usepackage[top=37mm,bottom=37mm,left=27mm,right=27mm]{geometry}
\sloppy
\pagestyle{myheadings}
\usepackage{amssymb}
\usepackage[latin1]{inputenc}
\usepackage{makeidx}
\makeindex
\usepackage{color}
\definecolor{FireBrick}{rgb}{0.5812,0.0074,0.0083}
\definecolor{RoyalBlue}{rgb}{0.0236,0.0894,0.6179}
\definecolor{RoyalGreen}{rgb}{0.0236,0.6179,0.0894}
\definecolor{RoyalRed}{rgb}{0.6179,0.0236,0.0894}
\definecolor{LightBlue}{rgb}{0.8544,0.9511,1.0000}
\definecolor{Black}{rgb}{0.0,0.0,0.0}

\definecolor{linkColor}{rgb}{0.0,0.0,0.554}
\definecolor{citeColor}{rgb}{0.0,0.0,0.554}
\definecolor{fileColor}{rgb}{0.0,0.0,0.554}
\definecolor{urlColor}{rgb}{0.0,0.0,0.554}
\definecolor{promptColor}{rgb}{0.0,0.0,0.589}
\definecolor{brkpromptColor}{rgb}{0.589,0.0,0.0}
\definecolor{gapinputColor}{rgb}{0.589,0.0,0.0}
\definecolor{gapoutputColor}{rgb}{0.0,0.0,0.0}

%%  for a long time these were red and blue by default,
%%  now black, but keep variables to overwrite
\definecolor{FuncColor}{rgb}{0.0,0.0,0.0}
%% strange name because of pdflatex bug:
\definecolor{Chapter }{rgb}{0.0,0.0,0.0}
\definecolor{DarkOlive}{rgb}{0.1047,0.2412,0.0064}


\usepackage{fancyvrb}

\usepackage{mathptmx,helvet}
\usepackage[T1]{fontenc}
\usepackage{textcomp}


\usepackage[
            pdftex=true,
            bookmarks=true,        
            a4paper=true,
            pdftitle={Written with GAPDoc},
            pdfcreator={LaTeX with hyperref package / GAPDoc},
            colorlinks=true,
            backref=page,
            breaklinks=true,
            linkcolor=linkColor,
            citecolor=citeColor,
            filecolor=fileColor,
            urlcolor=urlColor,
            pdfpagemode={UseNone}, 
           ]{hyperref}

\newcommand{\maintitlesize}{\fontsize{36}{38}\selectfont}

% write page numbers to a .pnr log file for online help
\newwrite\pagenrlog
\immediate\openout\pagenrlog =\jobname.pnr
\immediate\write\pagenrlog{PAGENRS := [}
\newcommand{\logpage}[1]{\protect\write\pagenrlog{#1, \thepage,}}
%% were never documented, give conflicts with some additional packages

\newcommand{\GAP}{\textsf{GAP}}

%% nicer description environments, allows long labels
\usepackage{enumitem}
\setdescription{style=nextline}

%% depth of toc
\setcounter{tocdepth}{1}






%% command for ColorPrompt style examples
\newcommand{\gapprompt}[1]{\color{promptColor}{\bfseries #1}}
\newcommand{\gapbrkprompt}[1]{\color{brkpromptColor}{\bfseries #1}}
\newcommand{\gapinput}[1]{\color{gapinputColor}{#1}}


\begin{document}

\logpage{[ 0, 0, 0 ]}
\begin{titlepage}
\mbox{}\vfill

\begin{center}{\maintitlesize \textbf{HPC\texttt{\symbol{45}}GAP \texttt{\symbol{45}}\texttt{\symbol{45}}\texttt{\symbol{45}} Reference Manual\\
\mbox{}}}\\
\vfill

\hypersetup{pdftitle={HPC\texttt{\symbol{45}}GAP \texttt{\symbol{45}}\texttt{\symbol{45}}\texttt{\symbol{45}} Reference Manual}}
\markright{\scriptsize \mbox{}\hfill HPC\texttt{\symbol{45}}GAP \texttt{\symbol{45}}\texttt{\symbol{45}}\texttt{\symbol{45}} Reference Manual \hfill\mbox{}}
{\Huge Release 4.17dev, today\mbox{}}\\[1cm]
\mbox{}\\[2cm]
{\Large \textbf{\strut The GAP Group    \strut\mbox{}}}\\
{\Large \textbf{\strut Reimer Behrends   \strut\mbox{}}}\\
{\Large \textbf{\strut Vladimir Janjic   \strut\mbox{}}}\\
\hypersetup{pdfauthor={The GAP Group   ; Reimer Behrends  ; Vladimir Janjic  }}
\end{center}\vfill

\mbox{}\\
{\mbox{}\\
\small \noindent \textbf{The GAP Group   }  Email: \href{mailto://support@gap-system.org} {\texttt{support@gap\texttt{\symbol{45}}system.org}}\\
  Homepage: \href{https://www.gap-system.org} {\texttt{https://www.gap\texttt{\symbol{45}}system.org}}}\\
{\mbox{}\\
\small \noindent \textbf{Reimer Behrends  }  Email: \href{mailto://behrends@gmail.com} {\texttt{behrends@gmail.com}}}\\
{\mbox{}\\
\small \noindent \textbf{Vladimir Janjic  }  Email: \href{mailto://vj32@st-andrews.ac.uk} {\texttt{vj32@st\texttt{\symbol{45}}andrews.ac.uk}}}\\
\end{titlepage}

\newpage\setcounter{page}{2}
{\small 
\section*{Copyright}
\logpage{[ 0, 0, 1 ]}
 Copyright {\copyright} (1987\texttt{\symbol{45}}this year) for the core part
of the \textsf{GAP} system by the \textsf{GAP} Group. 

 Most parts of this distribution, including the core part of the \textsf{GAP} system are distributed under the terms of the GNU General Public License, see \href{https://www.gnu.org/licenses/gpl.html} {\texttt{https://www.gnu.org/licenses/gpl.html}} or the file \texttt{GPL} in the \texttt{etc} directory of the \textsf{GAP} installation. 

 More detailed information about copyright and licenses of parts of this
distribution can be found in  (\textbf{Reference: Copyright and License}). 

 \textsf{GAP} is developed over a long time and has many authors and contributors. More
detailed information can be found in  (\textbf{Reference: Authors and Maintainers}). \mbox{}}\\[1cm]
\newpage

\def\contentsname{Contents\logpage{[ 0, 0, 2 ]}}

\tableofcontents
\newpage

 
\chapter{\textcolor{Chapter }{Tasks}}\label{Tasks}
\logpage{[ 1, 0, 0 ]}
\hyperdef{L}{X7EA285468533E153}{}
{
  
\section{\textcolor{Chapter }{Overview}}\label{Overview}
\logpage{[ 1, 1, 0 ]}
\hyperdef{L}{X8389AD927B74BA4A}{}
{
  Tasks provide mid\texttt{\symbol{45}} to high\texttt{\symbol{45}}level
functionality for programmers to describe asynchronous workflows. A task is an
asynchronously or synchronously executing job; functions exist to create tasks
that are executed concurrently, on demand, or in the current thread; to wait
for their completion, check their status, and retrieve any results. 

 Here is a simple example of sorting a list in the background: 
\begin{Verbatim}[commandchars=!@|,fontsize=\small,frame=single,label=Example]
  !gapprompt@gap>| !gapinput@task := RunTask(x -> SortedList(x), [3,2,1]);;|
  !gapprompt@gap>| !gapinput@WaitTask(task);|
  !gapprompt@gap>| !gapinput@TaskResult(task);|
  [ 1, 2, 3 ]
\end{Verbatim}
 \texttt{RunTask} (\ref{RunTask}) dispatches a task to run in the background; a task is described by a function
and zero or more arguments that are passed to \texttt{RunTask} (\ref{RunTask}). \texttt{WaitTask} (\ref{WaitTask}) waits for the task to complete; and \texttt{TaskResult} returns the result of the task. 

 \texttt{TaskResult} (\ref{TaskResult}) does an implicit \texttt{WaitTask} (\ref{WaitTask}), so the second line above can actually be omitted: 
\begin{Verbatim}[commandchars=!@|,fontsize=\small,frame=single,label=Example]
  !gapprompt@gap>| !gapinput@task := RunTask(x -> SortedList(x), [3,2,1]);;|
  !gapprompt@gap>| !gapinput@TaskResult(task);|
  [ 1, 2, 3 ]
\end{Verbatim}
 It is simple to run two tasks in parallel. Let's compute the factorial of
10000 by splitting the work between two tasks: 
\begin{Verbatim}[commandchars=!@|,fontsize=\small,frame=single,label=Example]
  !gapprompt@gap>| !gapinput@task1 := RunTask(Product, [1..5000]);;|
  !gapprompt@gap>| !gapinput@task2 := RunTask(Product, [5001..10000]);;|
  !gapprompt@gap>| !gapinput@TaskResult(task1) * TaskResult(task2) = Factorial(10000);|
  true
\end{Verbatim}
 You can use \texttt{DelayTask} (\ref{DelayTask}) to delay executing the task until its result is actually needed. 
\begin{Verbatim}[commandchars=!@|,fontsize=\small,frame=single,label=Example]
  !gapprompt@gap>| !gapinput@task1 := DelayTask(Product, [1..5000]);;|
  !gapprompt@gap>| !gapinput@task2 := DelayTask(Product, [5001..10000]);;|
  !gapprompt@gap>| !gapinput@WaitTask(task1, task2);|
  !gapprompt@gap>| !gapinput@TaskResult(task1) * TaskResult(task2) = Factorial(10000);|
  true
\end{Verbatim}
 Note that \texttt{WaitTask} (\ref{WaitTask}) is used here to start execution of both tasks; otherwise, \texttt{task2} would not be started until \texttt{TaskResult(task1)} has been evaluated. 

 To start execution of a delayed task, you can also use \texttt{ExecuteTask}. This has no effect if a task has already been running. 

 For convenience, you can also use \texttt{ImmediateTask} (\ref{ImmediateTask}) to execute a task synchronously (i.e., the task is started immediately and the
call does not return until the task has completed). 
\begin{Verbatim}[commandchars=!@|,fontsize=\small,frame=single,label=Example]
  !gapprompt@gap>| !gapinput@task := ImmediateTask(x -> SortedList(x), [3,2,1]);;|
  !gapprompt@gap>| !gapinput@TaskResult(task);|
  [ 1, 2, 3 ]
\end{Verbatim}
 This is indistinguishable from calling the function directly, but provides the
same interface as normal tasks. 

 If e.g. you want to call a function only for its
side\texttt{\symbol{45}}effects, it can be useful to ignore the result of a
task. \texttt{RunAsyncTask} (\ref{RunAsyncTask}) provides the necessary functionality. Such a task cannot be waited for and its
result (if any) is ignored. 
\begin{Verbatim}[commandchars=@|B,fontsize=\small,frame=single,label=Example]
  @gapprompt|gap>B @gapinput|RunAsyncTask(function() Print("Hello, world!\n"); end);;B
  @gapprompt|gap>B @gapinput|!listB
  --- Thread 0 [0]
  --- Thread 5 [5] (pending output)
  @gapprompt|gap>B @gapinput|!5B
  --- Switching to thread 5
  [5] Hello, world!
  !0
  --- Switching to thread 0
  gap>
\end{Verbatim}
 For more information on the multi\texttt{\symbol{45}}threaded user interface,
see Chapter \ref{Console User Interface}. 

 Task arguments are generally copied so that both the task that created them
and the task that uses them can access the data concurrently without fear of
race conditions. To avoid copying, arguments should be made shared or public
(see the relevant parts of section \ref{Migrating objects between regions} on migrating objects between regions); shared and public arguments will not be
copied. 

 HPC\texttt{\symbol{45}}GAP currently has multiple implementations of the task
API. To use an alternative implementation to the one documented here, set the
environment variable \texttt{GAP{\textunderscore}WORKSTEALING} to a non\texttt{\symbol{45}}empty value before starting GAP. }

 
\section{\textcolor{Chapter }{Running tasks}}\label{Running tasks}
\logpage{[ 1, 2, 0 ]}
\hyperdef{L}{X7C9F2EF384645B8F}{}
{
  

\subsection{\textcolor{Chapter }{RunTask}}
\logpage{[ 1, 2, 1 ]}\nobreak
\hyperdef{L}{X7DEC1649841E3813}{}
{\noindent\textcolor{FuncColor}{$\triangleright$\enspace\texttt{RunTask({\mdseries\slshape func[, arg1, ..., argn]})\index{RunTask@\texttt{RunTask}}
\label{RunTask}
}\hfill{\scriptsize (function)}}\\


 \texttt{RunTask} prepares a task for execution and starts it. The task will call the function \mbox{\texttt{\mdseries\slshape func}} with arguments \mbox{\texttt{\mdseries\slshape arg1}} through \mbox{\texttt{\mdseries\slshape argn}} (if provided). The return value of \mbox{\texttt{\mdseries\slshape func}} is the result of the task. The \texttt{RunTask} call itself returns a task object that can be used by functions that expect a
task argument. }

 

\subsection{\textcolor{Chapter }{ScheduleTask}}
\logpage{[ 1, 2, 2 ]}\nobreak
\hyperdef{L}{X7A86AE967D2A9E94}{}
{\noindent\textcolor{FuncColor}{$\triangleright$\enspace\texttt{ScheduleTask({\mdseries\slshape condition, func[, arg1, ..., argn]})\index{ScheduleTask@\texttt{ScheduleTask}}
\label{ScheduleTask}
}\hfill{\scriptsize (function)}}\\


 \texttt{ScheduleTask} prepares a task for execution, but, unlike \texttt{RunTask} (\ref{RunTask}) does not start it until \mbox{\texttt{\mdseries\slshape condition}} is met. See on how to construct conditions. Simple examples of conditions are
individual tasks, where execution occurs after the task completes, or lists of
tasks, where execution occurs after all tasks in the list complete. 
\begin{Verbatim}[commandchars=!@|,fontsize=\small,frame=single,label=Example]
  !gapprompt@gap>| !gapinput@t1 := RunTask(x->x*x, 3);;|
  !gapprompt@gap>| !gapinput@t2 := RunTask(x->x*x, 4);;|
  !gapprompt@gap>| !gapinput@t := ScheduleTask([t1, t2], function()|
  !gapprompt@>| !gapinput@          return TaskResult(t1) + TaskResult(t2);|
  !gapprompt@>| !gapinput@   end);;|
  !gapprompt@gap>| !gapinput@TaskResult(t);|
  25
\end{Verbatim}
 While the above example could also be achieved with \texttt{RunTask} (\ref{RunTask}) in lieu of \texttt{ScheduleTask}, since \texttt{TaskResult} (\ref{TaskResult}) would wait for \mbox{\texttt{\mdseries\slshape t1}} and \mbox{\texttt{\mdseries\slshape t2}} to complete, the above implementation does not actually start the final task
until the others are complete, making it more efficient, since no additional
worker thread needs to be occupied. }

 

\subsection{\textcolor{Chapter }{DelayTask}}
\logpage{[ 1, 2, 3 ]}\nobreak
\hyperdef{L}{X85BC2BBC7EF5B22E}{}
{\noindent\textcolor{FuncColor}{$\triangleright$\enspace\texttt{DelayTask({\mdseries\slshape func[, arg1, ..., argn]})\index{DelayTask@\texttt{DelayTask}}
\label{DelayTask}
}\hfill{\scriptsize (function)}}\\


 \texttt{DelayTask} works as \texttt{RunTask} (\ref{RunTask}), but its start is delayed until it is being waited for (including implicitly
by calling \texttt{TaskResult} (\ref{TaskResult})). }

 

\subsection{\textcolor{Chapter }{RunAsyncTask}}
\logpage{[ 1, 2, 4 ]}\nobreak
\hyperdef{L}{X856288C27C62B14A}{}
{\noindent\textcolor{FuncColor}{$\triangleright$\enspace\texttt{RunAsyncTask({\mdseries\slshape func[, arg1, ..., argn]})\index{RunAsyncTask@\texttt{RunAsyncTask}}
\label{RunAsyncTask}
}\hfill{\scriptsize (function)}}\\


 \texttt{RunAsyncTask} creates an asynchronous task. It works like \texttt{RunTask} (\ref{RunTask}), except that its result will be ignored. }

 

\subsection{\textcolor{Chapter }{ScheduleAsyncTask}}
\logpage{[ 1, 2, 5 ]}\nobreak
\hyperdef{L}{X7EAC7A4B86F5F74C}{}
{\noindent\textcolor{FuncColor}{$\triangleright$\enspace\texttt{ScheduleAsyncTask({\mdseries\slshape condition, func[, arg1, ..., argn]})\index{ScheduleAsyncTask@\texttt{ScheduleAsyncTask}}
\label{ScheduleAsyncTask}
}\hfill{\scriptsize (function)}}\\


 \texttt{ScheduleAsyncTask} creates an asynchronous task. It works like \texttt{ScheduleTask} (\ref{ScheduleTask}), except that its result will be ignored. }

 

\subsection{\textcolor{Chapter }{MakeTaskAsync}}
\logpage{[ 1, 2, 6 ]}\nobreak
\hyperdef{L}{X791CA0E07B16E7F9}{}
{\noindent\textcolor{FuncColor}{$\triangleright$\enspace\texttt{MakeTaskAsync({\mdseries\slshape task})\index{MakeTaskAsync@\texttt{MakeTaskAsync}}
\label{MakeTaskAsync}
}\hfill{\scriptsize (function)}}\\


 \texttt{MakeTaskAsync} turns a synchronous task into an asynchronous task that cannot be waited for
and whose result will be ignored. }

 

\subsection{\textcolor{Chapter }{ImmediateTask}}
\logpage{[ 1, 2, 7 ]}\nobreak
\hyperdef{L}{X84AFD0FA7A4E3633}{}
{\noindent\textcolor{FuncColor}{$\triangleright$\enspace\texttt{ImmediateTask({\mdseries\slshape func[, arg1, ..., argn]})\index{ImmediateTask@\texttt{ImmediateTask}}
\label{ImmediateTask}
}\hfill{\scriptsize (function)}}\\


 \texttt{ImmediateTask} executes the task specified by its arguments synchronously, usually within the
current thread. }

 

\subsection{\textcolor{Chapter }{ExecuteTask}}
\logpage{[ 1, 2, 8 ]}\nobreak
\hyperdef{L}{X83D2E32C8759FB37}{}
{\noindent\textcolor{FuncColor}{$\triangleright$\enspace\texttt{ExecuteTask({\mdseries\slshape task})\index{ExecuteTask@\texttt{ExecuteTask}}
\label{ExecuteTask}
}\hfill{\scriptsize (function)}}\\


 \texttt{ExecuteTask} starts \mbox{\texttt{\mdseries\slshape task}} if it is not already running. It has only an effect if its argument is a task
returned by \texttt{DelayTask} (\ref{DelayTask}); otherwise, it is a no\texttt{\symbol{45}}op. }

 

\subsection{\textcolor{Chapter }{WaitTask}}
\logpage{[ 1, 2, 9 ]}\nobreak
\hyperdef{L}{X861BC1488345C424}{}
{\noindent\textcolor{FuncColor}{$\triangleright$\enspace\texttt{WaitTask({\mdseries\slshape task1, ..., taskn})\index{WaitTask@\texttt{WaitTask}}
\label{WaitTask}
}\hfill{\scriptsize (function)}}\\
\noindent\textcolor{FuncColor}{$\triangleright$\enspace\texttt{WaitTask({\mdseries\slshape condition})\index{WaitTask@\texttt{WaitTask}!with a condition}
\label{WaitTask:with a condition}
}\hfill{\scriptsize (function)}}\\
\noindent\textcolor{FuncColor}{$\triangleright$\enspace\texttt{WaitTasks({\mdseries\slshape task1, ..., taskn})\index{WaitTasks@\texttt{WaitTasks}}
\label{WaitTasks}
}\hfill{\scriptsize (function)}}\\


 \texttt{WaitTask} waits until \mbox{\texttt{\mdseries\slshape task1}} through \mbox{\texttt{\mdseries\slshape taskn}} have completed; after that, it returns. Alternatively, a condition can be
passed to \texttt{WaitTask} in order to wait until a condition is met. See on how to construct conditions. \texttt{WaitTasks} is an alias for \texttt{WaitTask}. }

 

\subsection{\textcolor{Chapter }{WaitAnyTask}}
\logpage{[ 1, 2, 10 ]}\nobreak
\hyperdef{L}{X81489CD684EA11B6}{}
{\noindent\textcolor{FuncColor}{$\triangleright$\enspace\texttt{WaitAnyTask({\mdseries\slshape task1, ..., taskn})\index{WaitAnyTask@\texttt{WaitAnyTask}}
\label{WaitAnyTask}
}\hfill{\scriptsize (function)}}\\


 The \texttt{WaitAnyTask} function waits for any of its arguments to finish, then returns the number of
that task. 
\begin{Verbatim}[commandchars=!@|,fontsize=\small,frame=single,label=Example]
  !gapprompt@gap>| !gapinput@task1 := DelayTask(x->SortedList(x), [3,2,1]);;|
  !gapprompt@gap>| !gapinput@task2 := DelayTask(x->SortedList(x), [6,5,4]);;|
  !gapprompt@gap>| !gapinput@which := WaitAnyTask(task1, task2);|
  2
  !gapprompt@gap>| !gapinput@if which = 1 then|
  !gapprompt@>| !gapinput@     Display(TaskResult(task1));Display(TaskResult(task2));|
  !gapprompt@>| !gapinput@   else|
  !gapprompt@>| !gapinput@     Display(TaskResult(task2));Display(TaskResult(task1));|
  !gapprompt@>| !gapinput@   fi;|
  [ 4, 5, 6 ]
  [ 1, 2, 3 ]
\end{Verbatim}
 One can pass a list of tasks to \texttt{WaitAnyTask} as an argument; \texttt{WaitAnyTask([task1, ..., taskn])} behaves identically to \texttt{WaitAnyTask(task1, ..., taskn)}. }

 

\subsection{\textcolor{Chapter }{TaskResult}}
\logpage{[ 1, 2, 11 ]}\nobreak
\hyperdef{L}{X7A86C732823E04B6}{}
{\noindent\textcolor{FuncColor}{$\triangleright$\enspace\texttt{TaskResult({\mdseries\slshape task})\index{TaskResult@\texttt{TaskResult}}
\label{TaskResult}
}\hfill{\scriptsize (function)}}\\


 The \texttt{TaskResult} function returns the result of a task. It implicitly calls \texttt{WaitTask} (\ref{WaitTask}) if that is necessary. Multiple invocations of \texttt{TaskResult} with the same task argument will not do repeated waits and always return the
same value. 

 If the function executed by \mbox{\texttt{\mdseries\slshape task}} encounters an error, \texttt{TaskResult} returns \texttt{fail}. Whether \mbox{\texttt{\mdseries\slshape task}} encountered an error can be checked via \texttt{TaskSuccess} (\ref{TaskSuccess}). In case of an error, the error message can be retrieved via \texttt{TaskError} (\ref{TaskError}). }

 

\subsection{\textcolor{Chapter }{CullIdleTasks}}
\logpage{[ 1, 2, 12 ]}\nobreak
\hyperdef{L}{X830A82A7847993FF}{}
{\noindent\textcolor{FuncColor}{$\triangleright$\enspace\texttt{CullIdleTasks({\mdseries\slshape })\index{CullIdleTasks@\texttt{CullIdleTasks}}
\label{CullIdleTasks}
}\hfill{\scriptsize (function)}}\\


 This function terminates unused worker threads. }

 }

 
\section{\textcolor{Chapter }{Information about tasks}}\label{Information about tasks}
\logpage{[ 1, 3, 0 ]}
\hyperdef{L}{X86F8419B7968345B}{}
{
  

\subsection{\textcolor{Chapter }{TaskSuccess}}
\logpage{[ 1, 3, 1 ]}\nobreak
\hyperdef{L}{X87BB5EF38415B923}{}
{\noindent\textcolor{FuncColor}{$\triangleright$\enspace\texttt{TaskSuccess({\mdseries\slshape task})\index{TaskSuccess@\texttt{TaskSuccess}}
\label{TaskSuccess}
}\hfill{\scriptsize (function)}}\\


 \texttt{TaskSuccess} waits for \mbox{\texttt{\mdseries\slshape task}} and returns \texttt{true} if the it finished without encountering an error. Otherwise the function
returns \texttt{false}. }

 

\subsection{\textcolor{Chapter }{TaskError}}
\logpage{[ 1, 3, 2 ]}\nobreak
\hyperdef{L}{X85F42F5E7E04F3C2}{}
{\noindent\textcolor{FuncColor}{$\triangleright$\enspace\texttt{TaskError({\mdseries\slshape task})\index{TaskError@\texttt{TaskError}}
\label{TaskError}
}\hfill{\scriptsize (function)}}\\


 \texttt{TaskError} waits for \mbox{\texttt{\mdseries\slshape task}} and returns its error message, if it encountered an error. If it did not
encounter an error, the function returns \texttt{fail}. }

 

\subsection{\textcolor{Chapter }{CurrentTask}}
\logpage{[ 1, 3, 3 ]}\nobreak
\hyperdef{L}{X7C0C66E37DB9ACC0}{}
{\noindent\textcolor{FuncColor}{$\triangleright$\enspace\texttt{CurrentTask({\mdseries\slshape })\index{CurrentTask@\texttt{CurrentTask}}
\label{CurrentTask}
}\hfill{\scriptsize (function)}}\\


 The \texttt{CurrentTask} returns the currently running task. }

 

\subsection{\textcolor{Chapter }{RunningTasks}}
\logpage{[ 1, 3, 4 ]}\nobreak
\hyperdef{L}{X7DA4E06B792002C7}{}
{\noindent\textcolor{FuncColor}{$\triangleright$\enspace\texttt{RunningTasks({\mdseries\slshape })\index{RunningTasks@\texttt{RunningTasks}}
\label{RunningTasks}
}\hfill{\scriptsize (function)}}\\


 This function returns the number of currently running tasks. Note that it is
only an approximation and can change as new tasks are being started by other
threads. }

 

\subsection{\textcolor{Chapter }{TaskStarted}}
\logpage{[ 1, 3, 5 ]}\nobreak
\hyperdef{L}{X7A023A657BB4F17D}{}
{\noindent\textcolor{FuncColor}{$\triangleright$\enspace\texttt{TaskStarted({\mdseries\slshape task})\index{TaskStarted@\texttt{TaskStarted}}
\label{TaskStarted}
}\hfill{\scriptsize (function)}}\\


 This function returns true if the task has started executing (i.e., for any
non\texttt{\symbol{45}}delayed task), false otherwise. }

 

\subsection{\textcolor{Chapter }{TaskFinished}}
\logpage{[ 1, 3, 6 ]}\nobreak
\hyperdef{L}{X7E7500E17EC1067D}{}
{\noindent\textcolor{FuncColor}{$\triangleright$\enspace\texttt{TaskFinished({\mdseries\slshape task})\index{TaskFinished@\texttt{TaskFinished}}
\label{TaskFinished}
}\hfill{\scriptsize (function)}}\\


 This function returns true if the task has finished executing and its result
is available, false otherwise. }

 

\subsection{\textcolor{Chapter }{TaskIsAsync}}
\logpage{[ 1, 3, 7 ]}\nobreak
\hyperdef{L}{X7B37CED386024A1F}{}
{\noindent\textcolor{FuncColor}{$\triangleright$\enspace\texttt{TaskIsAsync({\mdseries\slshape task})\index{TaskIsAsync@\texttt{TaskIsAsync}}
\label{TaskIsAsync}
}\hfill{\scriptsize (function)}}\\


 This function returns true if the task is asynchronous, true otherwise. }

 }

 
\section{\textcolor{Chapter }{Cancelling tasks}}\label{Cancelling tasks}
\logpage{[ 1, 4, 0 ]}
\hyperdef{L}{X7B286945821A3BBA}{}
{
  HPC\texttt{\symbol{45}}GAP uses a cooperative model for task cancellation. A
programmer can request the cancellation of another task, but it is up to that
other task to actually terminate itself. The tasks library has functions to
request cancellation, to test for the cancellation state of a task, and to
perform actions in response to cancellation requests. 

\subsection{\textcolor{Chapter }{CancelTask}}
\logpage{[ 1, 4, 1 ]}\nobreak
\hyperdef{L}{X86C26E3081C3C9AA}{}
{\noindent\textcolor{FuncColor}{$\triangleright$\enspace\texttt{CancelTask({\mdseries\slshape task})\index{CancelTask@\texttt{CancelTask}}
\label{CancelTask}
}\hfill{\scriptsize (function)}}\\


 \texttt{CancelTask} submits a request that \mbox{\texttt{\mdseries\slshape task}} is to be cancelled. }

 

\subsection{\textcolor{Chapter }{TaskCancellationRequested}}
\logpage{[ 1, 4, 2 ]}\nobreak
\hyperdef{L}{X8431B91487A8A3F1}{}
{\noindent\textcolor{FuncColor}{$\triangleright$\enspace\texttt{TaskCancellationRequested({\mdseries\slshape task})\index{TaskCancellationRequested@\texttt{TaskCancellationRequested}}
\label{TaskCancellationRequested}
}\hfill{\scriptsize (function)}}\\


 \texttt{TaskCancellationRequested} returns true if \texttt{CancelTask} (\ref{CancelTask}) has been called for \mbox{\texttt{\mdseries\slshape task}}, false otherwise. }

 

\subsection{\textcolor{Chapter }{OnTaskCancellation}}
\logpage{[ 1, 4, 3 ]}\nobreak
\hyperdef{L}{X7D61F9C487232DE7}{}
{\noindent\textcolor{FuncColor}{$\triangleright$\enspace\texttt{OnTaskCancellation({\mdseries\slshape exit{\textunderscore}func})\index{OnTaskCancellation@\texttt{OnTaskCancellation}}
\label{OnTaskCancellation}
}\hfill{\scriptsize (function)}}\\


 \texttt{OnTaskCancellation} tests if cancellation for the current task has been requested. If so, then \mbox{\texttt{\mdseries\slshape exit{\textunderscore}func}} will be called (as a parameterless function) and the current task will be
aborted. The result of the current task will be the value of \mbox{\texttt{\mdseries\slshape exit{\textunderscore}func()}}. 
\begin{Verbatim}[commandchars=!@|,fontsize=\small,frame=single,label=Example]
  !gapprompt@gap>| !gapinput@task := RunTask(function()|
  !gapprompt@>| !gapinput@     while true do|
  !gapprompt@>| !gapinput@       OnTaskCancellation(function() return 314; end);|
  !gapprompt@>| !gapinput@     od;|
  !gapprompt@>| !gapinput@   end);;|
  !gapprompt@gap>| !gapinput@CancelTask(task);|
  !gapprompt@gap>| !gapinput@TaskResult(task);|
  314
\end{Verbatim}
 }

 

\subsection{\textcolor{Chapter }{OnTaskCancellationReturn}}
\logpage{[ 1, 4, 4 ]}\nobreak
\hyperdef{L}{X82D1773F8058CA54}{}
{\noindent\textcolor{FuncColor}{$\triangleright$\enspace\texttt{OnTaskCancellationReturn({\mdseries\slshape value})\index{OnTaskCancellationReturn@\texttt{OnTaskCancellationReturn}}
\label{OnTaskCancellationReturn}
}\hfill{\scriptsize (function)}}\\


 \texttt{OnTaskCancellationReturn} is a convenience function that does the same as: \texttt{OnTaskCancellation(function() return value; end);} }

 }

 
\section{\textcolor{Chapter }{Conditions}}\label{Conditions}
\logpage{[ 1, 5, 0 ]}
\hyperdef{L}{X802FC4C284F68696}{}
{
  \texttt{ScheduleTask} (\ref{ScheduleTask}) and \texttt{WaitTask} (\ref{WaitTask}) can be made to wait on more complex conditions than just tasks. A condition is
either a milestone, a task, or a list of milestones and tasks. \texttt{ScheduleTask} (\ref{ScheduleTask}) starts its task and \texttt{WaitTask} (\ref{WaitTask}) returns when the condition has been met. A condition represented by a task is
met when the task has completed. A condition represented by a milestone is met
when the milestone has been achieved (see below). A condition represented by a
list is met when all conditions in the list have been met. }

 
\section{\textcolor{Chapter }{Milestones}}\label{Milestones}
\logpage{[ 1, 6, 0 ]}
\hyperdef{L}{X7EE6DC40842CD74B}{}
{
  Milestones are a way to represent abstract conditions to which multiple tasks
can contribute. 

\subsection{\textcolor{Chapter }{NewMilestone}}
\logpage{[ 1, 6, 1 ]}\nobreak
\hyperdef{L}{X80CB3247854A46E4}{}
{\noindent\textcolor{FuncColor}{$\triangleright$\enspace\texttt{NewMilestone({\mdseries\slshape [list]})\index{NewMilestone@\texttt{NewMilestone}}
\label{NewMilestone}
}\hfill{\scriptsize (function)}}\\


 The \texttt{NewMilestone} function creates a new milestone. Its argument is a list of targets, which
must be a list of integers and/or strings. If omitted, the list defaults to \texttt{[0]}. }

 

\subsection{\textcolor{Chapter }{ContributeToMilestone}}
\logpage{[ 1, 6, 2 ]}\nobreak
\hyperdef{L}{X800AE2737CB715C5}{}
{\noindent\textcolor{FuncColor}{$\triangleright$\enspace\texttt{ContributeToMilestone({\mdseries\slshape milestone, target})\index{ContributeToMilestone@\texttt{ContributeToMilestone}}
\label{ContributeToMilestone}
}\hfill{\scriptsize (function)}}\\


 The \texttt{ContributeToMilestone} milestone function contributes the specified target to the milestone. Once all
targets have been contributed to a milestone, it has been achieved. }

 

\subsection{\textcolor{Chapter }{AchieveMilestone}}
\logpage{[ 1, 6, 3 ]}\nobreak
\hyperdef{L}{X7ACB26D38096CD56}{}
{\noindent\textcolor{FuncColor}{$\triangleright$\enspace\texttt{AchieveMilestone({\mdseries\slshape milestone})\index{AchieveMilestone@\texttt{AchieveMilestone}}
\label{AchieveMilestone}
}\hfill{\scriptsize (function)}}\\


 The \texttt{AchieveMilestone} function allows a program to achieve a milestone in a single step without
adding individual targets to it. This is most useful in conjunction with the
default value for \texttt{NewMilestone} (\ref{NewMilestone}), e.g. 
\begin{Verbatim}[commandchars=!@|,fontsize=\small,frame=single,label=Example]
  !gapprompt@gap>| !gapinput@m := NewMilestone();;|
  !gapprompt@gap>| !gapinput@AchieveMilestone(m);|
\end{Verbatim}
 }

{\textgreater} 

\subsection{\textcolor{Chapter }{IsMilestoneAchieved}}
\logpage{[ 1, 6, 4 ]}\nobreak
\hyperdef{L}{X83E7AF007E040FB9}{}
{\noindent\textcolor{FuncColor}{$\triangleright$\enspace\texttt{IsMilestoneAchieved({\mdseries\slshape milestone})\index{IsMilestoneAchieved@\texttt{IsMilestoneAchieved}}
\label{IsMilestoneAchieved}
}\hfill{\scriptsize (function)}}\\


 \texttt{IsMilestoneAchieved} tests explicitly if a milestone has been achieved. It returns \texttt{true} on success, \texttt{false} otherwise. 
\begin{Verbatim}[commandchars=!@|,fontsize=\small,frame=single,label=Example]
  !gapprompt@gap>| !gapinput@m := NewMilestone([1,2]);;|
  !gapprompt@gap>| !gapinput@ContributeToMilestone(m, 1);|
  !gapprompt@gap>| !gapinput@IsMilestoneAchieved(m);|
  false
  !gapprompt@gap>| !gapinput@ContributeToMilestone(m, 2);|
  !gapprompt@gap>| !gapinput@IsMilestoneAchieved(m);|
  true
\end{Verbatim}
 }

 }

 }

 
\chapter{\textcolor{Chapter }{Variables in HPC\texttt{\symbol{45}}GAP}}\label{Variables in HPC-GAP}
\logpage{[ 2, 0, 0 ]}
\hyperdef{L}{X8000D4CD7F4F5594}{}
{
  Variables with global scope have revised semantics in
HPC\texttt{\symbol{45}}GAP in order to address concurrency issues. The normal
semantics of global variables that are only accessed by a single thread remain
unaltered. 
\section{\textcolor{Chapter }{Global variables}}\label{Global variables}
\logpage{[ 2, 1, 0 ]}
\hyperdef{L}{X7D9044767BEB1523}{}
{
  Global variables in HPC\texttt{\symbol{45}}GAP can be accessed by all threads
concurrently without explicit synchronization. Concurrent access is safe, but
it is not deterministic. If multiple threads attempt to modify the same global
variable simultaneously, the resulting value of the variable is random; it
will be one of the values assigned by a thread, but it is impossible to
predict with certainty which specific one will be assigned. }

 
\section{\textcolor{Chapter }{Thread\texttt{\symbol{45}}local variables}}\label{Thread-local variables}
\logpage{[ 2, 2, 0 ]}
\hyperdef{L}{X7D93681D7B5E8DCD}{}
{
  HPC\texttt{\symbol{45}}GAP supports the notion of
thread\texttt{\symbol{45}}local variables. Thread\texttt{\symbol{45}}local
variables are (after being declared as such) accessed and modified like global
variables. However, unlike global variables, each thread can assign a distinct
value to a thread\texttt{\symbol{45}}local variable. 
\begin{Verbatim}[commandchars=!@|,fontsize=\small,frame=single,label=Example]
  !gapprompt@gap>| !gapinput@MakeThreadLocal("x");|
  !gapprompt@gap>| !gapinput@x := 1;;|
  !gapprompt@gap>| !gapinput@WaitTask(RunTask(function() x := 2; end));|
  !gapprompt@gap>| !gapinput@x;|
  1
\end{Verbatim}
 As can be seen here, the assignment to \texttt{x} in a separate thread does not overwrite the value of \texttt{x} in the main thread. 

\subsection{\textcolor{Chapter }{MakeThreadLocal}}
\logpage{[ 2, 2, 1 ]}\nobreak
\hyperdef{L}{X7FE1310180B55506}{}
{\noindent\textcolor{FuncColor}{$\triangleright$\enspace\texttt{MakeThreadLocal({\mdseries\slshape name})\index{MakeThreadLocal@\texttt{MakeThreadLocal}}
\label{MakeThreadLocal}
}\hfill{\scriptsize (function)}}\\


 \texttt{MakeThreadLocal} makes the variable described by the string \mbox{\texttt{\mdseries\slshape name}} a thread\texttt{\symbol{45}}local variable. It normally does not give it an
initial value; either explicit per\texttt{\symbol{45}}thread assignment or a
call to \texttt{BindThreadLocal} (\ref{BindThreadLocal}) or \texttt{BindThreadLocalConstructor} (\ref{BindThreadLocalConstructor}) to provide a default value is necessary. 

 If a global variable with the same name exists and is bound at the time of the
call, its value will be used as the default value as though \texttt{BindThreadLocal} (\ref{BindThreadLocal}) had been called with that value as its second argument. }

 

\subsection{\textcolor{Chapter }{BindThreadLocal}}
\logpage{[ 2, 2, 2 ]}\nobreak
\hyperdef{L}{X81F4832C7ED44627}{}
{\noindent\textcolor{FuncColor}{$\triangleright$\enspace\texttt{BindThreadLocal({\mdseries\slshape name, obj})\index{BindThreadLocal@\texttt{BindThreadLocal}}
\label{BindThreadLocal}
}\hfill{\scriptsize (function)}}\\


 \texttt{BindThreadLocal} gives the thread\texttt{\symbol{45}}local variable described by the string \mbox{\texttt{\mdseries\slshape name}} the default value \mbox{\texttt{\mdseries\slshape obj}}. The first time the thread\texttt{\symbol{45}}local variable is accessed in a
thread thereafter, it will yield \mbox{\texttt{\mdseries\slshape obj}} as its value if it hasn't been assigned a specific value yet. }

 

\subsection{\textcolor{Chapter }{BindThreadLocalConstructor}}
\logpage{[ 2, 2, 3 ]}\nobreak
\hyperdef{L}{X8606D69B82B8AE84}{}
{\noindent\textcolor{FuncColor}{$\triangleright$\enspace\texttt{BindThreadLocalConstructor({\mdseries\slshape name, func})\index{BindThreadLocalConstructor@\texttt{BindThreadLocalConstructor}}
\label{BindThreadLocalConstructor}
}\hfill{\scriptsize (function)}}\\


 \texttt{BindThreadLocal} (\ref{BindThreadLocal}) gives the thread\texttt{\symbol{45}}local variable described by the string \mbox{\texttt{\mdseries\slshape name}} the constructor \mbox{\texttt{\mdseries\slshape func}}. The first time the thread\texttt{\symbol{45}}local variable is accessed in a
thread thereafter, it will yield \mbox{\texttt{\mdseries\slshape func()}} as its value if it hasn't been assigned a specific value yet. }

 

\subsection{\textcolor{Chapter }{ThreadVar}}
\logpage{[ 2, 2, 4 ]}\nobreak
\hyperdef{L}{X7C508AFD8115AEB0}{}
{\noindent\textcolor{FuncColor}{$\triangleright$\enspace\texttt{ThreadVar\index{ThreadVar@\texttt{ThreadVar}}
\label{ThreadVar}
}\hfill{\scriptsize (global variable)}}\\


 All thread\texttt{\symbol{45}}local variables are stored in the
thread\texttt{\symbol{45}}local record \texttt{ThreadVar}. Thus, if \texttt{x} is a thread\texttt{\symbol{45}}local variable, using \texttt{ThreadVar.x} is the same as using \texttt{x}. }

 }

 }

 
\chapter{\textcolor{Chapter }{How HPC\texttt{\symbol{45}}GAP organizes shared memory: Regions}}\label{How HPC-GAP organizes shared memory: Regions}
\logpage{[ 3, 0, 0 ]}
\hyperdef{L}{X8076353A830B11B6}{}
{
  HPC\texttt{\symbol{45}}GAP allows multiple threads to access data shared
between them; to avoid common concurrency errors, such as race conditions, it
partitions GAP objects into regions. Access to regions is regulated so that no
two threads can modify objects in the same region at the same time and so that
objects that are being read by one thread cannot concurrently be modified by
another. 
\section{\textcolor{Chapter }{Thread\texttt{\symbol{45}}local regions}}\label{Thread-local regions}
\logpage{[ 3, 1, 0 ]}
\hyperdef{L}{X7E2199568017C74F}{}
{
  Each thread has an associated thread\texttt{\symbol{45}}local region. When a
thread implicitly or explicitly creates a new object, that object initially
belongs to the thread's thread\texttt{\symbol{45}}local region. 

 Only the thread can read or modify objects in its
thread\texttt{\symbol{45}}local region. For other threads to access an object,
that object has to be migrated into a different region first. }

 
\section{\textcolor{Chapter }{Shared regions}}\label{Shared regions}
\logpage{[ 3, 2, 0 ]}
\hyperdef{L}{X7B697BA17A813E7D}{}
{
  Shared regions are explicitly created through the \texttt{ShareObj} (\ref{ShareObj}) and \texttt{ShareSingleObj} (\ref{ShareSingleObj}) primitives (see below). Multiple threads can access them concurrently, but
accessing them requires that a thread uses an \texttt{atomic} statement to acquire a read or write lock beforehand. 

 See the section on \texttt{atomic} statements (\ref{The atomic statement.}) for details. }

 
\section{\textcolor{Chapter }{Ordering of shared regions}}\label{Ordering of shared regions}
\logpage{[ 3, 3, 0 ]}
\hyperdef{L}{X83627591876D3FF3}{}
{
  Shared regions are by default ordered; each shared region has an associated
numeric precedence level. Regions can generally only be locked in order of
descending precedence. The purpose of this mechanism is to avoid accidental
deadlocks. 

 The ordering requirement can be overridden in two ways: regions with a
negative precedence are excluded from it. This exception should be used with
care, as it can lead to deadlocks. 

 Alternatively, two or more regions can be locked simultaneously via the \texttt{atomic} statement. In this case, the ordering of these regions relative to each other
can be arbitrary. }

 
\section{\textcolor{Chapter }{The public region}}\label{The public region}
\logpage{[ 3, 4, 0 ]}
\hyperdef{L}{X8239FDC583A4E39D}{}
{
  A special public region contains objects that only permit atomic operations.
These include, in particular, all immutable objects (immutable in the sense
that their in\texttt{\symbol{45}}memory representation cannot change). 

 All threads can access objects in the public region at all times without
needing to acquire a read\texttt{\symbol{45}} or write\texttt{\symbol{45}}lock
beforehand. }

 
\section{\textcolor{Chapter }{The read\texttt{\symbol{45}}only region}}\label{The read-only region}
\logpage{[ 3, 5, 0 ]}
\hyperdef{L}{X7E0116957AFB982D}{}
{
  The read\texttt{\symbol{45}}only region is another special region that
contains objects that are only meant to be read; attempting to modify an
object in that region will result in a runtime error. To obtain a modifiable
copy of such an object, the \texttt{CopyRegion} (\ref{CopyRegion}) primitive can be used. }

 
\section{\textcolor{Chapter }{Migrating objects between regions}}\label{Migrating objects between regions}
\logpage{[ 3, 6, 0 ]}
\hyperdef{L}{X87315FA584A37637}{}
{
  Objects can be migrated between regions using a number of functions. In order
to migrate an object, the current thread must have exclusive access to that
object; the object must be in its thread\texttt{\symbol{45}}local region or it
must be in a shared region for which the current thread holds a write lock. 

 The \texttt{ShareObj} (\ref{ShareObj}) and \texttt{ShareSingleObj} (\ref{ShareSingleObj}) functions create a new shared region and migrate their respective argument to
that region; \texttt{ShareObj} will also migrate all subobjects that are within the same region, while \texttt{ShareSingleObj} will leave the subobjects unaffected. 

 The \texttt{MigrateObj} (\ref{MigrateObj}) and \texttt{MigrateSingleObj} (\ref{MigrateSingleObj}) functions migrate objects to an existing region. The first argument of either
function is the object to be migrated; the second is either a region (as
returned by the \texttt{RegionOf} (\ref{RegionOf}) function) or an object whose containing region the first argument is to be
migrated to. 

 The current thread needs exclusive access to the target region (denoted by the
second argument) for the operation to succeed. If successful, the first
argument will be in the same region as the second argument afterwards. In the
case of \texttt{MigrateObj} (\ref{MigrateObj}), all subobjects within the same region as the first argument will also be
migrated to the target region. 

 Finally, \texttt{AdoptObj} (\ref{AdoptObj}) and \texttt{AdoptSingleObj} (\ref{AdoptSingleObj}) are special cases of \texttt{MigrateObj} (\ref{MigrateObj}) and \texttt{MigrateSingleObj} (\ref{MigrateSingleObj}), where the target region is the thread\texttt{\symbol{45}}local region of the
current thread. 

 To migrate objects to the read\texttt{\symbol{45}}only region, one can use \texttt{MakeReadOnlyObj} (\ref{MakeReadOnlyObj}) and \texttt{MakeReadOnlySingleObj} (\ref{MakeReadOnlySingleObj}). The first migrates its argument and all its subjobjects that are within the
same region to the read\texttt{\symbol{45}}only region; the second migrates
only the argument itself, but not its subobjects. 

 It is generally not possible to migrate objects explicitly to the public
region; only objects with purely atomic operations can be made public and that
is done automatically when they are created. 

 The exception are immutable objects. When \texttt{MakeImmutable} (\textbf{Reference: MakeImmutable}) is used, its argument is automatically moved to the public region. 
\begin{Verbatim}[commandchars=!@|,fontsize=\small,frame=single,label=Example]
  !gapprompt@gap>| !gapinput@RegionOf(MakeImmutable([1,2,3]));|
  <public region>
\end{Verbatim}
 }

 
\section{\textcolor{Chapter }{Region names}}\label{Region names}
\logpage{[ 3, 7, 0 ]}
\hyperdef{L}{X79A959BF7C24234F}{}
{
  Regions can be given names, either explicitly via \texttt{SetRegionName} (\ref{SetRegionName}) or when they are created via \texttt{ShareObj} (\ref{ShareObj}) and \texttt{ShareSingleObj} (\ref{ShareSingleObj}). Thread\texttt{\symbol{45}}local regions, the public, and the readonly region
are given names by default. 

 Multiple regions can have the same name. }

 
\section{\textcolor{Chapter }{Controlling access to regions}}\label{Controlling access to regions}
\logpage{[ 3, 8, 0 ]}
\hyperdef{L}{X827637EE7A69AFCD}{}
{
  If either GAP code or a kernel primitive attempts to access an object that it
is not allowed to access according to these semantics, either a "write guard
error" (for a failed write access) or a "read guard error" (for a failed read
access) will be raised. The global variable \texttt{LastInaccessible} will contain the object that caused such an error. 

 One exception is that threads can modify objects in regions that they have
only read access (but not write access) to using write\texttt{\symbol{45}}once
functions \texttt{\symbol{45}} see section \ref{Write-once functionality}. 

 To inspect objects whose contents lie in other regions (and therefore cannot
be displayed by \texttt{PrintObj} (\textbf{Reference: PrintObj}) or \texttt{ViewObj} (\textbf{Reference: ViewObj}), the functions \texttt{ViewShared} (\ref{ViewShared}) and \texttt{UNSAFE{\textunderscore}VIEW} (\ref{UNSAFEuScoreVIEW}) can be used. }

 
\section{\textcolor{Chapter }{Functions relating to regions}}\label{Functions relating to regions}
\logpage{[ 3, 9, 0 ]}
\hyperdef{L}{X7BE832987B1DC975}{}
{
  

\subsection{\textcolor{Chapter }{NewRegion}}
\logpage{[ 3, 9, 1 ]}\nobreak
\hyperdef{L}{X851C5F3C82F6F5AE}{}
{\noindent\textcolor{FuncColor}{$\triangleright$\enspace\texttt{NewRegion({\mdseries\slshape [name, ]prec})\index{NewRegion@\texttt{NewRegion}}
\label{NewRegion}
}\hfill{\scriptsize (function)}}\\


 The function \texttt{NewRegion} creates a new shared region. If the optional argument \mbox{\texttt{\mdseries\slshape name}} is provided, then the name of the new region will be set to \mbox{\texttt{\mdseries\slshape name}}. 
\begin{Verbatim}[commandchars=!@|,fontsize=\small,frame=single,label=Example]
  !gapprompt@gap>| !gapinput@NewRegion("example region");|
  <region: example region>
\end{Verbatim}
 \texttt{NewRegion} will create a region with a high precedence level. It is intended to be called
by user code. The exact precedence level can be adjusted with \mbox{\texttt{\mdseries\slshape prec}}, which must be an integer in the range \texttt{[\texttt{\symbol{45}}1000..1000]}; \mbox{\texttt{\mdseries\slshape prec}} will be added to the normal precedence level. }

 

\subsection{\textcolor{Chapter }{NewLibraryRegion}}
\logpage{[ 3, 9, 2 ]}\nobreak
\hyperdef{L}{X83864D427DE991F2}{}
{\noindent\textcolor{FuncColor}{$\triangleright$\enspace\texttt{NewLibraryRegion({\mdseries\slshape [name, ]prec})\index{NewLibraryRegion@\texttt{NewLibraryRegion}}
\label{NewLibraryRegion}
}\hfill{\scriptsize (function)}}\\


 \texttt{NewLibraryRegion} functions like \texttt{NewRegion} (\ref{NewRegion}), except that the precedence of the region it creates is below that of \texttt{NewRegion} (\ref{NewRegion}). It is intended to be used by user libraries and GAP packages. }

 

\subsection{\textcolor{Chapter }{NewSystemRegion}}
\logpage{[ 3, 9, 3 ]}\nobreak
\hyperdef{L}{X7FB0BE4C78CA85DA}{}
{\noindent\textcolor{FuncColor}{$\triangleright$\enspace\texttt{NewSystemRegion({\mdseries\slshape [name, ]prec})\index{NewSystemRegion@\texttt{NewSystemRegion}}
\label{NewSystemRegion}
}\hfill{\scriptsize (function)}}\\


 \texttt{NewSystemRegion} functions like \texttt{NewRegion} (\ref{NewRegion}), except that the precedence of the region it creates is below that of \texttt{NewLibraryRegion} (\ref{NewLibraryRegion}). It is intended to be used by the standard GAP library. }

 

\subsection{\textcolor{Chapter }{NewKernelRegion}}
\logpage{[ 3, 9, 4 ]}\nobreak
\hyperdef{L}{X825A881A7A39C5C3}{}
{\noindent\textcolor{FuncColor}{$\triangleright$\enspace\texttt{NewKernelRegion({\mdseries\slshape [name, ]prec})\index{NewKernelRegion@\texttt{NewKernelRegion}}
\label{NewKernelRegion}
}\hfill{\scriptsize (function)}}\\


 \texttt{NewKernelRegion} functions like \texttt{NewRegion} (\ref{NewRegion}), except that the precedence of the region it creates is below that of \texttt{NewSystemRegion} (\ref{NewSystemRegion}). It is intended to be used by the GAP kernel, and GAP library code that
interacts closely with the kernel. }

 

\subsection{\textcolor{Chapter }{NewInternalRegion}}
\logpage{[ 3, 9, 5 ]}\nobreak
\hyperdef{L}{X86C54C9278FE00F4}{}
{\noindent\textcolor{FuncColor}{$\triangleright$\enspace\texttt{NewInternalRegion({\mdseries\slshape [name]})\index{NewInternalRegion@\texttt{NewInternalRegion}}
\label{NewInternalRegion}
}\hfill{\scriptsize (function)}}\\


 \texttt{NewInternalRegion} functions like \texttt{NewRegion} (\ref{NewRegion}), except that the precedence of the region it creates is the lowest available.
It is intended to be used for regions that are
self\texttt{\symbol{45}}contained; i.e. no function that uses such a region
may lock another region while accessing it. The precedence level of an
internal region cannot be adjusted. }

 

\subsection{\textcolor{Chapter }{NewSpecialRegion}}
\logpage{[ 3, 9, 6 ]}\nobreak
\hyperdef{L}{X7A7FFA847E090257}{}
{\noindent\textcolor{FuncColor}{$\triangleright$\enspace\texttt{NewSpecialRegion({\mdseries\slshape [name]})\index{NewSpecialRegion@\texttt{NewSpecialRegion}}
\label{NewSpecialRegion}
}\hfill{\scriptsize (function)}}\\


 \texttt{NewLibraryRegion} (\ref{NewLibraryRegion}) functions like \texttt{NewRegion} (\ref{NewRegion}), except that the precedence of the region it creates is negative. It is thus
exempt from normal ordering and deadlock checks. }

 

\subsection{\textcolor{Chapter }{RegionOf}}
\logpage{[ 3, 9, 7 ]}\nobreak
\hyperdef{L}{X86BEBBAF855AA26A}{}
{\noindent\textcolor{FuncColor}{$\triangleright$\enspace\texttt{RegionOf({\mdseries\slshape obj})\index{RegionOf@\texttt{RegionOf}}
\label{RegionOf}
}\hfill{\scriptsize (function)}}\\


 
\begin{Verbatim}[commandchars=!@|,fontsize=\small,frame=single,label=Example]
  !gapprompt@gap>| !gapinput@RegionOf(1/2);|
  <public region>
  !gapprompt@gap>| !gapinput@RegionOf([1,2,3]);|
  <region: thread region #0>
  !gapprompt@gap>| !gapinput@RegionOf(ShareObj([1,2,3]));|
  <region 0x45deaa0>
  !gapprompt@gap>| !gapinput@RegionOf(ShareObj([1,2,3]));|
  <region 0x45deaa0>
  !gapprompt@gap>| !gapinput@RegionOf(ShareObj([1,2,3], "test region"));|
  <region: test region>
\end{Verbatim}
 Note that the unique number that each region is identified with is
system\texttt{\symbol{45}}specific and can change each time the code is being
run. Region objects returned by \texttt{RegionOf} can be compared: 
\begin{Verbatim}[commandchars=!@|,fontsize=\small,frame=single,label=Example]
  !gapprompt@gap>| !gapinput@RegionOf([1,2,3]) = RegionOf([4,5,6]);|
  true
\end{Verbatim}
 The result in this example is true because both lists are in the same
thread\texttt{\symbol{45}}local region. }

 

\subsection{\textcolor{Chapter }{RegionPrecedence}}
\logpage{[ 3, 9, 8 ]}\nobreak
\hyperdef{L}{X87421870782B33C7}{}
{\noindent\textcolor{FuncColor}{$\triangleright$\enspace\texttt{RegionPrecedence({\mdseries\slshape obj})\index{RegionPrecedence@\texttt{RegionPrecedence}}
\label{RegionPrecedence}
}\hfill{\scriptsize (function)}}\\


 \texttt{RegionPrecedence} will return the precedence of the region of \mbox{\texttt{\mdseries\slshape obj}}. 
\begin{Verbatim}[commandchars=!@|,fontsize=\small,frame=single,label=Example]
  !gapprompt@gap>| !gapinput@RegionPrecedence(NewRegion("Test"));|
  30000
  !gapprompt@gap>| !gapinput@RegionPrecedence(NewRegion("Test2", 1));|
  30001
  !gapprompt@gap>| !gapinput@RegionPrecedence(NewLibraryRegion("LibTest", -1));|
  19999
\end{Verbatim}
 }

 

\subsection{\textcolor{Chapter }{ShareObj}}
\logpage{[ 3, 9, 9 ]}\nobreak
\hyperdef{L}{X7D5982617A3027BD}{}
{\noindent\textcolor{FuncColor}{$\triangleright$\enspace\texttt{ShareObj({\mdseries\slshape obj[[, name], prec]})\index{ShareObj@\texttt{ShareObj}}
\label{ShareObj}
}\hfill{\scriptsize (function)}}\\


 The \texttt{ShareObj} function creates a new shared region and migrates the object and all its
subobjects to that region. If the optional argument \mbox{\texttt{\mdseries\slshape name}} is provided, then the name of the new region is set to \mbox{\texttt{\mdseries\slshape name}}. 

 \texttt{ShareObj} will create a region with a high precedence level. It is intended to be called
by user code. The actual precedence level can be adjusted by the optional \mbox{\texttt{\mdseries\slshape prec}} argument in the same way as for \texttt{NewRegion} (\ref{NewRegion}). }

 

\subsection{\textcolor{Chapter }{ShareLibraryObj}}
\logpage{[ 3, 9, 10 ]}\nobreak
\hyperdef{L}{X79E455D27E12C5B4}{}
{\noindent\textcolor{FuncColor}{$\triangleright$\enspace\texttt{ShareLibraryObj({\mdseries\slshape obj[[, name], prec]})\index{ShareLibraryObj@\texttt{ShareLibraryObj}}
\label{ShareLibraryObj}
}\hfill{\scriptsize (function)}}\\


 \texttt{ShareLibraryObj} functions like \texttt{ShareObj} (\ref{ShareObj}), except that the precedence of the region it creates is below that of \texttt{ShareObj} (\ref{ShareObj}). It is intended to be used by user libraries and GAP packages. }

 

\subsection{\textcolor{Chapter }{ShareSystemObj}}
\logpage{[ 3, 9, 11 ]}\nobreak
\hyperdef{L}{X867DE843791EEF65}{}
{\noindent\textcolor{FuncColor}{$\triangleright$\enspace\texttt{ShareSystemObj({\mdseries\slshape obj[[, name], prec]})\index{ShareSystemObj@\texttt{ShareSystemObj}}
\label{ShareSystemObj}
}\hfill{\scriptsize (function)}}\\


 \texttt{ShareSystemObj} functions like \texttt{ShareObj} (\ref{ShareObj}), except that the precedence of the region it creates is below that of \texttt{ShareLibraryObj} (\ref{ShareLibraryObj}). It is intended to be used by the standard GAP library. }

 

\subsection{\textcolor{Chapter }{ShareKernelObj}}
\logpage{[ 3, 9, 12 ]}\nobreak
\hyperdef{L}{X7E7540D17EADF97A}{}
{\noindent\textcolor{FuncColor}{$\triangleright$\enspace\texttt{ShareKernelObj({\mdseries\slshape obj[[, name], prec]})\index{ShareKernelObj@\texttt{ShareKernelObj}}
\label{ShareKernelObj}
}\hfill{\scriptsize (function)}}\\


 \texttt{ShareKernelObj} functions like \texttt{ShareObj} (\ref{ShareObj}), except that the precedence of the region it creates is below that of \texttt{ShareSystemObj} (\ref{ShareSystemObj}). It is intended to be used by the GAP kernel, and GAP library code that
interacts closely with the kernel. }

 

\subsection{\textcolor{Chapter }{ShareInternalObj}}
\logpage{[ 3, 9, 13 ]}\nobreak
\hyperdef{L}{X792DAE2C83BD1554}{}
{\noindent\textcolor{FuncColor}{$\triangleright$\enspace\texttt{ShareInternalObj({\mdseries\slshape obj[, name]})\index{ShareInternalObj@\texttt{ShareInternalObj}}
\label{ShareInternalObj}
}\hfill{\scriptsize (function)}}\\


 \texttt{ShareInternalObj} functions like \texttt{ShareObj} (\ref{ShareObj}), except that the precedence of the region it creates is the lowest available.
It is intended to be used for regions that are
self\texttt{\symbol{45}}contained; i.e. no function that uses such a region
may lock another region while accessing it. }

 

\subsection{\textcolor{Chapter }{ShareSpecialObj}}
\logpage{[ 3, 9, 14 ]}\nobreak
\hyperdef{L}{X82F3B2597E0EC15E}{}
{\noindent\textcolor{FuncColor}{$\triangleright$\enspace\texttt{ShareSpecialObj({\mdseries\slshape obj[, name]})\index{ShareSpecialObj@\texttt{ShareSpecialObj}}
\label{ShareSpecialObj}
}\hfill{\scriptsize (function)}}\\


 \texttt{ShareSpecialObj} functions like \texttt{ShareObj} (\ref{ShareObj}), except that the precedence of the region it creates is negative. It is thus
exempt from normal ordering and deadlock checks. }

 

\subsection{\textcolor{Chapter }{ShareSingleObj}}
\logpage{[ 3, 9, 15 ]}\nobreak
\hyperdef{L}{X8508A72B7C215FA5}{}
{\noindent\textcolor{FuncColor}{$\triangleright$\enspace\texttt{ShareSingleObj({\mdseries\slshape obj[[, name], prec]})\index{ShareSingleObj@\texttt{ShareSingleObj}}
\label{ShareSingleObj}
}\hfill{\scriptsize (function)}}\\


 The \texttt{ShareSingleObj} function creates a new shared region and migrates the object, but not its
subobjects, to that region. If the optional argument \mbox{\texttt{\mdseries\slshape name}} is provided, then the name of the new region is set to \mbox{\texttt{\mdseries\slshape name}}. 
\begin{Verbatim}[commandchars=!@|,fontsize=\small,frame=single,label=Example]
  !gapprompt@gap>| !gapinput@m := [ [1, 2], [3, 4] ];;|
  !gapprompt@gap>| !gapinput@ShareSingleObj(m);;|
  !gapprompt@gap>| !gapinput@atomic readonly m do|
  !gapprompt@>| !gapinput@     Display([ IsShared(m), IsShared(m[1]), IsShared(m[2]) ]);|
  !gapprompt@>| !gapinput@   od;|
  [ true, false, false ]
\end{Verbatim}
 \texttt{ShareSingleObj} will create a region with a high precedence level. It is intended to be called
by user code. The actual precedence level can be adjusted by the optional \mbox{\texttt{\mdseries\slshape prec}} argument in the same way as for \texttt{NewRegion} (\ref{NewRegion}). }

 

\subsection{\textcolor{Chapter }{ShareSingleLibraryObj}}
\logpage{[ 3, 9, 16 ]}\nobreak
\hyperdef{L}{X87A1962578CDA61D}{}
{\noindent\textcolor{FuncColor}{$\triangleright$\enspace\texttt{ShareSingleLibraryObj({\mdseries\slshape obj[[, name], prec]})\index{ShareSingleLibraryObj@\texttt{ShareSingleLibraryObj}}
\label{ShareSingleLibraryObj}
}\hfill{\scriptsize (function)}}\\


 \texttt{ShareSingleLibraryObj} functions like \texttt{ShareSingleObj} (\ref{ShareSingleObj}), except that the precedence of the region it creates is below that of \texttt{ShareSingleObj} (\ref{ShareSingleObj}). It is intended to be used by user libraries and GAP packages. }

 

\subsection{\textcolor{Chapter }{ShareSingleSystemObj}}
\logpage{[ 3, 9, 17 ]}\nobreak
\hyperdef{L}{X8352EF8B83390656}{}
{\noindent\textcolor{FuncColor}{$\triangleright$\enspace\texttt{ShareSingleSystemObj({\mdseries\slshape obj[[, name], prec]})\index{ShareSingleSystemObj@\texttt{ShareSingleSystemObj}}
\label{ShareSingleSystemObj}
}\hfill{\scriptsize (function)}}\\


 \texttt{ShareSingleSystemObj} functions like \texttt{ShareSingleObj} (\ref{ShareSingleObj}), except that the precedence of the region it creates is below that of \texttt{ShareSingleLibraryObj} (\ref{ShareSingleLibraryObj}). It is intended to be used by the standard GAP library. }

 

\subsection{\textcolor{Chapter }{ShareSingleKernelObj}}
\logpage{[ 3, 9, 18 ]}\nobreak
\hyperdef{L}{X7B5A471982EFD292}{}
{\noindent\textcolor{FuncColor}{$\triangleright$\enspace\texttt{ShareSingleKernelObj({\mdseries\slshape obj[[, name], prec]})\index{ShareSingleKernelObj@\texttt{ShareSingleKernelObj}}
\label{ShareSingleKernelObj}
}\hfill{\scriptsize (function)}}\\


 \texttt{ShareSingleKernelObj} functions like \texttt{ShareSingleObj} (\ref{ShareSingleObj}), except that the precedence of the region it creates is below that of \texttt{ShareSingleSystemObj} (\ref{ShareSingleSystemObj}). It is intended to be used by the GAP kernel, and GAP library code that
interacts closely with the kernel. }

 

\subsection{\textcolor{Chapter }{ShareSingleInternalObj}}
\logpage{[ 3, 9, 19 ]}\nobreak
\hyperdef{L}{X85C5F5A67DAFD919}{}
{\noindent\textcolor{FuncColor}{$\triangleright$\enspace\texttt{ShareSingleInternalObj({\mdseries\slshape obj[, name]})\index{ShareSingleInternalObj@\texttt{ShareSingleInternalObj}}
\label{ShareSingleInternalObj}
}\hfill{\scriptsize (function)}}\\


 \texttt{ShareSingleInternalObj} functions like \texttt{ShareSingleObj} (\ref{ShareSingleObj}), except that the precedence of the region it creates is the lowest available.
It is intended to be used for regions that are
self\texttt{\symbol{45}}contained; i.e. no function that uses such a region
may lock another region while accessing it. }

 

\subsection{\textcolor{Chapter }{ShareSingleSpecialObj}}
\logpage{[ 3, 9, 20 ]}\nobreak
\hyperdef{L}{X7CB671AE7A411314}{}
{\noindent\textcolor{FuncColor}{$\triangleright$\enspace\texttt{ShareSingleSpecialObj({\mdseries\slshape obj[, name]})\index{ShareSingleSpecialObj@\texttt{ShareSingleSpecialObj}}
\label{ShareSingleSpecialObj}
}\hfill{\scriptsize (function)}}\\


 \texttt{ShareSingleLibraryObj} (\ref{ShareSingleLibraryObj}) functions like \texttt{ShareSingleObj} (\ref{ShareSingleObj}), except that the precedence of the region it creates is negative. It is thus
exempt from normal ordering and deadlock checks. }

 

\subsection{\textcolor{Chapter }{MigrateObj}}
\logpage{[ 3, 9, 21 ]}\nobreak
\hyperdef{L}{X81A356DD84E76A8A}{}
{\noindent\textcolor{FuncColor}{$\triangleright$\enspace\texttt{MigrateObj({\mdseries\slshape obj, target})\index{MigrateObj@\texttt{MigrateObj}}
\label{MigrateObj}
}\hfill{\scriptsize (function)}}\\


 The \texttt{MigrateObj} function migrates \mbox{\texttt{\mdseries\slshape obj}} (and all subobjects contained within the same region) to the region denoted by
the \mbox{\texttt{\mdseries\slshape target}} argument. Here, \mbox{\texttt{\mdseries\slshape target}} can either be a region object returned by \texttt{RegionOf} (\ref{RegionOf}) or a normal gap object. If \mbox{\texttt{\mdseries\slshape target}} is a normal gap object, \mbox{\texttt{\mdseries\slshape obj}} will be migrated to the region containing \texttt{target}. 

 For the operation to succeed, the current thread must have exclusive access to
the target region and the object being migrated. }

 

\subsection{\textcolor{Chapter }{MigrateSingleObj}}
\logpage{[ 3, 9, 22 ]}\nobreak
\hyperdef{L}{X7BAE5A7282793684}{}
{\noindent\textcolor{FuncColor}{$\triangleright$\enspace\texttt{MigrateSingleObj({\mdseries\slshape obj, target})\index{MigrateSingleObj@\texttt{MigrateSingleObj}}
\label{MigrateSingleObj}
}\hfill{\scriptsize (function)}}\\


 The \texttt{MigrateSingleObj} function works like \texttt{MigrateObj} (\ref{MigrateObj}), except that it does not migrate the subobjects of \texttt{obj}. }

 

\subsection{\textcolor{Chapter }{LockAndMigrateObj}}
\logpage{[ 3, 9, 23 ]}\nobreak
\hyperdef{L}{X7E4B54BF837E81C0}{}
{\noindent\textcolor{FuncColor}{$\triangleright$\enspace\texttt{LockAndMigrateObj({\mdseries\slshape obj, target})\index{LockAndMigrateObj@\texttt{LockAndMigrateObj}}
\label{LockAndMigrateObj}
}\hfill{\scriptsize (function)}}\\


 The \texttt{LockAndMigrateObj} function works like \texttt{MigrateObj} (\ref{MigrateObj}), except that it will automatically try to acquire a lock for the region
containing \mbox{\texttt{\mdseries\slshape target}} if it does not have one already. }

 

\subsection{\textcolor{Chapter }{IncorporateObj}}
\logpage{[ 3, 9, 24 ]}\nobreak
\hyperdef{L}{X7D1943AF793296F7}{}
{\noindent\textcolor{FuncColor}{$\triangleright$\enspace\texttt{IncorporateObj({\mdseries\slshape target, index, value})\index{IncorporateObj@\texttt{IncorporateObj}}
\label{IncorporateObj}
}\hfill{\scriptsize (function)}}\\


 The \texttt{IncorporateObj} function allows convenient migration to a shared list or record. If \mbox{\texttt{\mdseries\slshape target}} is a list, then \texttt{IncorporateObj} is equivalent to: 
\begin{Verbatim}[commandchars=!@|,fontsize=\small,frame=single,label=Example]
  IncorporateObj := function(target, index, value)
    atomic value do
      target[index] := MigrateObj(value, target)
    od;
  end;
\end{Verbatim}
 If \mbox{\texttt{\mdseries\slshape target}} is a record, then it is equivalent to: 
\begin{Verbatim}[commandchars=!@|,fontsize=\small,frame=single,label=Example]
  IncorporateObj := function(target, index, value)
    atomic value do
      target.(index) := MigrateObj(value, target)
    od;
  end;
\end{Verbatim}
 The intended purpose is the population of a shared list or record with values
after its creation. Example: 
\begin{Verbatim}[commandchars=!@|,fontsize=\small,frame=single,label=Example]
  !gapprompt@gap>| !gapinput@list := ShareObj([]);|
  !gapprompt@gap>| !gapinput@atomic list do|
  !gapprompt@>| !gapinput@     IncorporateObj(list, 1, [1,2,3]);|
  !gapprompt@>| !gapinput@     IncorporateObj(list, 2, [4,5,6]);|
  !gapprompt@>| !gapinput@     IncorporateObj(list, 3, [7,8,9]);|
  !gapprompt@>| !gapinput@   od;|
  !gapprompt@gap>| !gapinput@ViewShared(list);|
  [ [ 1, 2, 3 ], [ 4, 5, 6 ], [ 7, 8, 9 ] ]
\end{Verbatim}
 Using plain assignment would leave the newly created lists in the
thread\texttt{\symbol{45}}local region. }

 

\subsection{\textcolor{Chapter }{AtomicIncorporateObj}}
\logpage{[ 3, 9, 25 ]}\nobreak
\hyperdef{L}{X876843717F4437CB}{}
{\noindent\textcolor{FuncColor}{$\triangleright$\enspace\texttt{AtomicIncorporateObj({\mdseries\slshape target, index, value})\index{AtomicIncorporateObj@\texttt{AtomicIncorporateObj}}
\label{AtomicIncorporateObj}
}\hfill{\scriptsize (function)}}\\


 \texttt{AtomicIncorporateObj} extends \texttt{IncorporateObj} (\ref{IncorporateObj}) by also locking the target. I.e., for a list, it is equivalent to: 
\begin{Verbatim}[commandchars=!@|,fontsize=\small,frame=single,label=Example]
  AtomicIncorporateObj := function(target, index, value)
    atomic target, value do
      target[index] := MigrateObj(value, target)
    od;
  end;
\end{Verbatim}
 If \texttt{target} is a record, then it is equivalent to: 
\begin{Verbatim}[commandchars=!@|,fontsize=\small,frame=single,label=Example]
  AtomicIncorporateObj := function(target, index, value)
    atomic value do
      target.(index) := MigrateObj(value, target)
    od;
  end;
\end{Verbatim}
 }

 

\subsection{\textcolor{Chapter }{AdoptObj}}
\logpage{[ 3, 9, 26 ]}\nobreak
\hyperdef{L}{X784C978D801191E2}{}
{\noindent\textcolor{FuncColor}{$\triangleright$\enspace\texttt{AdoptObj({\mdseries\slshape obj})\index{AdoptObj@\texttt{AdoptObj}}
\label{AdoptObj}
}\hfill{\scriptsize (function)}}\\


 The \texttt{AdoptObj} function migrates \mbox{\texttt{\mdseries\slshape obj}} (and all its subobjects contained within the same region) to the thread's
current region. It requires exclusive access to \mbox{\texttt{\mdseries\slshape obj}}. 
\begin{Verbatim}[commandchars=!@|,fontsize=\small,frame=single,label=Example]
  !gapprompt@gap>| !gapinput@l := ShareObj([1,2,3]);;|
  !gapprompt@gap>| !gapinput@IsThreadLocal(l);|
  false
  !gapprompt@gap>| !gapinput@atomic l do AdoptObj(l); od;|
  !gapprompt@gap>| !gapinput@IsThreadLocal(l);|
  true
\end{Verbatim}
 }

 

\subsection{\textcolor{Chapter }{AdoptSingleObj}}
\logpage{[ 3, 9, 27 ]}\nobreak
\hyperdef{L}{X834DDB388600E9FA}{}
{\noindent\textcolor{FuncColor}{$\triangleright$\enspace\texttt{AdoptSingleObj({\mdseries\slshape obj})\index{AdoptSingleObj@\texttt{AdoptSingleObj}}
\label{AdoptSingleObj}
}\hfill{\scriptsize (function)}}\\


 The \texttt{AdoptSingleObj} function works like \texttt{AdoptObj} (\ref{AdoptObj}), except that it does not migrate the subobjects of \mbox{\texttt{\mdseries\slshape obj}}. }

 

\subsection{\textcolor{Chapter }{LockAndAdoptObj}}
\logpage{[ 3, 9, 28 ]}\nobreak
\hyperdef{L}{X867CDC9285D30DE8}{}
{\noindent\textcolor{FuncColor}{$\triangleright$\enspace\texttt{LockAndAdoptObj({\mdseries\slshape obj})\index{LockAndAdoptObj@\texttt{LockAndAdoptObj}}
\label{LockAndAdoptObj}
}\hfill{\scriptsize (function)}}\\


 The \texttt{LockAndAdoptObj} function works like \texttt{AdoptObj} (\ref{AdoptObj}), except that it will attempt acquire an exclusive lock for the region
containing \mbox{\texttt{\mdseries\slshape obj}} if it does not have one already. }

 

\subsection{\textcolor{Chapter }{CopyRegion}}
\logpage{[ 3, 9, 29 ]}\nobreak
\hyperdef{L}{X7C71A88487762733}{}
{\noindent\textcolor{FuncColor}{$\triangleright$\enspace\texttt{CopyRegion({\mdseries\slshape obj})\index{CopyRegion@\texttt{CopyRegion}}
\label{CopyRegion}
}\hfill{\scriptsize (function)}}\\


 The \texttt{CopyRegion} function performs a structural copy of \mbox{\texttt{\mdseries\slshape obj}}. The resulting objects will be located in the current thread's
thread\texttt{\symbol{45}}local region. The function returns the copy as its
result. 
\begin{Verbatim}[commandchars=!@|,fontsize=\small,frame=single,label=Example]
  !gapprompt@gap>| !gapinput@l := MakeReadOnlyObj([1,2,3]);|
  [ 1, 2, 3 ]
  !gapprompt@gap>| !gapinput@l2 := CopyRegion(l);|
  [ 1, 2, 3 ]
  !gapprompt@gap>| !gapinput@RegionOf(l) = RegionOf(l2);|
  false
  !gapprompt@gap>| !gapinput@IsIdenticalObj(l, l2);|
  false
  !gapprompt@gap>| !gapinput@l = l2;|
  true
\end{Verbatim}
 }

 

\subsection{\textcolor{Chapter }{IsPublic}}
\logpage{[ 3, 9, 30 ]}\nobreak
\hyperdef{L}{X8222929685E9959A}{}
{\noindent\textcolor{FuncColor}{$\triangleright$\enspace\texttt{IsPublic({\mdseries\slshape obj})\index{IsPublic@\texttt{IsPublic}}
\label{IsPublic}
}\hfill{\scriptsize (function)}}\\


 The \texttt{IsPublic} function returns true if its argument is an object in the public region, false
otherwise. 
\begin{Verbatim}[commandchars=!@|,fontsize=\small,frame=single,label=Example]
  !gapprompt@gap>| !gapinput@IsPublic(1/2);|
  true
  !gapprompt@gap>| !gapinput@IsPublic([1,2,3]);|
  false
  !gapprompt@gap>| !gapinput@IsPublic(ShareObj([1,2,3]));|
  false
  !gapprompt@gap>| !gapinput@IsPublic(MakeImmutable([1,2,3]));|
  true
\end{Verbatim}
 }

 

\subsection{\textcolor{Chapter }{IsThreadLocal}}
\logpage{[ 3, 9, 31 ]}\nobreak
\hyperdef{L}{X86B2EEF67C3378F0}{}
{\noindent\textcolor{FuncColor}{$\triangleright$\enspace\texttt{IsThreadLocal({\mdseries\slshape obj})\index{IsThreadLocal@\texttt{IsThreadLocal}}
\label{IsThreadLocal}
}\hfill{\scriptsize (function)}}\\


 The \texttt{IsThreadLocal} function returns true if its argument is an object in the current thread's
thread\texttt{\symbol{45}}local region, false otherwise. 
\begin{Verbatim}[commandchars=!@|,fontsize=\small,frame=single,label=Example]
  !gapprompt@gap>| !gapinput@IsThreadLocal([1,2,3]);|
  true
  !gapprompt@gap>| !gapinput@IsThreadLocal(ShareObj([1,2,3]));|
  false
  !gapprompt@gap>| !gapinput@IsThreadLocal(1/2);|
  false
  !gapprompt@gap>| !gapinput@RegionOf(1/2);|
  <public region>
\end{Verbatim}
 }

 

\subsection{\textcolor{Chapter }{IsShared}}
\logpage{[ 3, 9, 32 ]}\nobreak
\hyperdef{L}{X80A11F3C84DB512E}{}
{\noindent\textcolor{FuncColor}{$\triangleright$\enspace\texttt{IsShared({\mdseries\slshape obj})\index{IsShared@\texttt{IsShared}}
\label{IsShared}
}\hfill{\scriptsize (function)}}\\


 The \texttt{IsShared} function returns true if its argument is an object in a shared region. Note
that if the current thread does not hold a lock on that shared region, another
thread can migrate \mbox{\texttt{\mdseries\slshape obj}} to a different region before the result is being evaluated; this can lead to
race conditions. The function is intended primarily for debugging, not to
build actual program logic around. }

 

\subsection{\textcolor{Chapter }{HaveReadAccess}}
\logpage{[ 3, 9, 33 ]}\nobreak
\hyperdef{L}{X827A26A67C99316C}{}
{\noindent\textcolor{FuncColor}{$\triangleright$\enspace\texttt{HaveReadAccess({\mdseries\slshape obj})\index{HaveReadAccess@\texttt{HaveReadAccess}}
\label{HaveReadAccess}
}\hfill{\scriptsize (function)}}\\


 The \texttt{HaveReadAccess} function returns true if the current thread has read access to \mbox{\texttt{\mdseries\slshape obj}}. 
\begin{Verbatim}[commandchars=!@|,fontsize=\small,frame=single,label=Example]
  !gapprompt@gap>| !gapinput@HaveReadAccess([1,2,3]);|
  true
  !gapprompt@gap>| !gapinput@l := ShareObj([1,2,3]);;|
  !gapprompt@gap>| !gapinput@HaveReadAccess(l);|
  false
  !gapprompt@gap>| !gapinput@atomic readonly l do t := HaveReadAccess(l); od;; t;|
  true
\end{Verbatim}
 }

 

\subsection{\textcolor{Chapter }{HaveWriteAccess}}
\logpage{[ 3, 9, 34 ]}\nobreak
\hyperdef{L}{X794206E5845006EA}{}
{\noindent\textcolor{FuncColor}{$\triangleright$\enspace\texttt{HaveWriteAccess({\mdseries\slshape obj})\index{HaveWriteAccess@\texttt{HaveWriteAccess}}
\label{HaveWriteAccess}
}\hfill{\scriptsize (function)}}\\


 The \texttt{HaveWriteAccess} function returns true if the current thread has write access to \mbox{\texttt{\mdseries\slshape obj}}. 
\begin{Verbatim}[commandchars=!@|,fontsize=\small,frame=single,label=Example]
  !gapprompt@gap>| !gapinput@HaveWriteAccess([1,2,3]);|
  true
  !gapprompt@gap>| !gapinput@l := ShareObj([1,2,3]);;|
  !gapprompt@gap>| !gapinput@HaveWriteAccess(l);|
  false
  !gapprompt@gap>| !gapinput@atomic readwrite l do t := HaveWriteAccess(l); od;; t;|
  true
\end{Verbatim}
 }

 

\subsection{\textcolor{Chapter }{MakeReadOnlyObj}}
\logpage{[ 3, 9, 35 ]}\nobreak
\hyperdef{L}{X7F53D70285AF37B4}{}
{\noindent\textcolor{FuncColor}{$\triangleright$\enspace\texttt{MakeReadOnlyObj({\mdseries\slshape obj})\index{MakeReadOnlyObj@\texttt{MakeReadOnlyObj}}
\label{MakeReadOnlyObj}
}\hfill{\scriptsize (function)}}\\


 The \texttt{MakeReadOnlyObj} function migrates \mbox{\texttt{\mdseries\slshape obj}} and all its subobjects that are within the same region as \mbox{\texttt{\mdseries\slshape obj}} to the read\texttt{\symbol{45}}only region. It returns \mbox{\texttt{\mdseries\slshape obj}}. }

 

\subsection{\textcolor{Chapter }{MakeReadOnlySingleObj}}
\logpage{[ 3, 9, 36 ]}\nobreak
\hyperdef{L}{X7EC9341A865BCC35}{}
{\noindent\textcolor{FuncColor}{$\triangleright$\enspace\texttt{MakeReadOnlySingleObj({\mdseries\slshape obj})\index{MakeReadOnlySingleObj@\texttt{MakeReadOnlySingleObj}}
\label{MakeReadOnlySingleObj}
}\hfill{\scriptsize (function)}}\\


 The \texttt{MakeReadOnlySingleObj} function migrates \mbox{\texttt{\mdseries\slshape obj}}, but not any of its subobjects, to the read\texttt{\symbol{45}}only region.
It returns \mbox{\texttt{\mdseries\slshape obj}}. }

 

\subsection{\textcolor{Chapter }{IsReadOnlyObj}}
\logpage{[ 3, 9, 37 ]}\nobreak
\hyperdef{L}{X860008F57EFE21C4}{}
{\noindent\textcolor{FuncColor}{$\triangleright$\enspace\texttt{IsReadOnlyObj({\mdseries\slshape obj})\index{IsReadOnlyObj@\texttt{IsReadOnlyObj}}
\label{IsReadOnlyObj}
}\hfill{\scriptsize (function)}}\\


 The \texttt{IsReadOnlyObj} function returns \texttt{true} if \mbox{\texttt{\mdseries\slshape obj}} is in the read\texttt{\symbol{45}}only region, \texttt{false} otherwise. 
\begin{Verbatim}[commandchars=!@|,fontsize=\small,frame=single,label=Example]
  !gapprompt@gap>| !gapinput@IsReadOnlyObj([1,2,3]);|
  false
  !gapprompt@gap>| !gapinput@IsReadOnlyObj(MakeImmutable([1,2,3]));|
  false
  !gapprompt@gap>| !gapinput@IsReadOnlyObj(MakeReadOnlyObj([1,2,3]));|
  true
\end{Verbatim}
 }

 

\subsection{\textcolor{Chapter }{SetRegionName}}
\logpage{[ 3, 9, 38 ]}\nobreak
\hyperdef{L}{X7F1E2F707F72371E}{}
{\noindent\textcolor{FuncColor}{$\triangleright$\enspace\texttt{SetRegionName({\mdseries\slshape obj, name})\index{SetRegionName@\texttt{SetRegionName}}
\label{SetRegionName}
}\hfill{\scriptsize (function)}}\\


 The \texttt{SetRegionName} function sets the name of the region of \mbox{\texttt{\mdseries\slshape obj}} to \mbox{\texttt{\mdseries\slshape name}}. }

 

\subsection{\textcolor{Chapter }{ClearRegionName}}
\logpage{[ 3, 9, 39 ]}\nobreak
\hyperdef{L}{X8427E1537ADC4575}{}
{\noindent\textcolor{FuncColor}{$\triangleright$\enspace\texttt{ClearRegionName({\mdseries\slshape obj})\index{ClearRegionName@\texttt{ClearRegionName}}
\label{ClearRegionName}
}\hfill{\scriptsize (function)}}\\


 The \texttt{ClearRegionName} function clears the name of the region of \mbox{\texttt{\mdseries\slshape obj}} to \mbox{\texttt{\mdseries\slshape name}}. }

 

\subsection{\textcolor{Chapter }{RegionName}}
\logpage{[ 3, 9, 40 ]}\nobreak
\hyperdef{L}{X7959FC997CC9177C}{}
{\noindent\textcolor{FuncColor}{$\triangleright$\enspace\texttt{RegionName({\mdseries\slshape obj})\index{RegionName@\texttt{RegionName}}
\label{RegionName}
}\hfill{\scriptsize (function)}}\\


 The \texttt{RegionName} function returns the name of the region of \mbox{\texttt{\mdseries\slshape obj}}. If that region does not have a name, \texttt{fail} will be returned. }

 

\subsection{\textcolor{Chapter }{ViewShared}}
\logpage{[ 3, 9, 41 ]}\nobreak
\hyperdef{L}{X80D0DFAB7F7241E8}{}
{\noindent\textcolor{FuncColor}{$\triangleright$\enspace\texttt{ViewShared({\mdseries\slshape obj})\index{ViewShared@\texttt{ViewShared}}
\label{ViewShared}
}\hfill{\scriptsize (function)}}\\


 The \texttt{ViewShared} function allows the inspection of objects in shared regions. It will try to
lock the region and then call \texttt{ViewObj(obj)}. If it cannot acquire a lock for the region, it will simply display the
normal description of the object. }

 

\subsection{\textcolor{Chapter }{UNSAFE{\textunderscore}VIEW}}
\logpage{[ 3, 9, 42 ]}\nobreak
\hyperdef{L}{X7FD39BCC8526AC53}{}
{\noindent\textcolor{FuncColor}{$\triangleright$\enspace\texttt{UNSAFE{\textunderscore}VIEW({\mdseries\slshape obj})\index{UNSAFE{\textunderscore}VIEW@\texttt{UNSAFE{\textunderscore}VIEW}}
\label{UNSAFEuScoreVIEW}
}\hfill{\scriptsize (function)}}\\


 The \texttt{UNSAFE{\textunderscore}VIEW} function allows the inspection of any object in the system, regardless of
whether the current thread has access to the region containing it. It should
be used with care: If the object inspected is being modified by another thread
concurrently, the resulting behavior is undefined. 

 Moreover, the function works by temporarily disabling read and write guards
for regions, so other threads may corrupt memory rather than producing errors. 

 It is generally safe to use if all threads but the current one are paused. }

 
\subsection{\textcolor{Chapter }{The \texttt{atomic} statement.}}\label{The atomic statement.}
\logpage{[ 3, 9, 43 ]}
\hyperdef{L}{X7FD1B1B785E24734}{}
{
  \index{atomic@\texttt{atomic}!no value} The \texttt{atomic} statement ensures exclusive or read\texttt{\symbol{45}}only access to one or
more shared regions for statements within its scope. It has the following
syntax: 
\begin{Verbatim}[commandchars=!@A,fontsize=\small,frame=single,label=Example]
  atomic ([readwrite|readonly] expr (, expr)* )* do
    statements
  od;
\end{Verbatim}
 Each expression is evaluated and the region containing the resulting object is
locked with either a read\texttt{\symbol{45}}write or
read\texttt{\symbol{45}}only lock, depending on the keyword preceding the
expression. If neither the \texttt{readwrite} nor the \texttt{readonly} keyword was provided, read\texttt{\symbol{45}}write locks are used by default.
Examples: 
\begin{Verbatim}[commandchars=!@|,fontsize=\small,frame=single,label=Example]
  !gapprompt@gap>| !gapinput@l := ShareObj([1,2,3]);;|
  !gapprompt@gap>| !gapinput@atomic readwrite l do l[3] := 9; od;|
  !gapprompt@gap>| !gapinput@atomic l do l[2] := 4; od;|
  !gapprompt@gap>| !gapinput@atomic readonly l do Display(l); od;|
  [ 1, 4, 9 ]
\end{Verbatim}
 
\begin{Verbatim}[commandchars=!@|,fontsize=\small,frame=single,label=Example]
  !gapprompt@gap>| !gapinput@l := ShareObj([1,2,3,4,5]);;|
  !gapprompt@gap>| !gapinput@l2 := ShareObj([6,7,8]);;|
  !gapprompt@gap>| !gapinput@atomic readwrite l, readonly l2 do|
  !gapprompt@>| !gapinput@     for i in [1..3] do l[i] := l2[i]; od;|
  !gapprompt@>| !gapinput@     l3 := AdoptObj(l);|
  !gapprompt@>| !gapinput@   od;|
  !gapprompt@gap>| !gapinput@l3;|
  [ 6, 7, 8, 4, 5 ]
\end{Verbatim}
 Atomic statements must observe region ordering. That means that the highest
precedence level of a region locked by an atomic statement must be less than
the lowest precedene level of a region that is locked by the same thread at
the time the atomic statement is executed. }

 }

 
\section{\textcolor{Chapter }{Atomic functions}}\label{Atomic functions}
\logpage{[ 3, 10, 0 ]}
\hyperdef{L}{X8789D7A57CFC13BC}{}
{
  Instead of atomic regions, entire functions can be declared to be atomic. This
has the same effect as though the function's body were enclosed in an atomic
statement. Function arguments can be declared either \texttt{readwrite} or \texttt{readonly}; they will be locked in the same way as for a lock statement. If a function
argument is preceded by neither \texttt{readwrite} nor \texttt{readonly}, the corresponding object will not be locked. Example: 
\begin{Verbatim}[commandchars=!@|,fontsize=\small,frame=single,label=Example]
  !gapprompt@gap>| !gapinput@AddAtomic := atomic function(readwrite list, readonly item)|
  !gapprompt@>| !gapinput@     Add(list, item);|
  !gapprompt@>| !gapinput@   end;|
\end{Verbatim}
 }

 
\section{\textcolor{Chapter }{Write\texttt{\symbol{45}}once functionality}}\label{Write-once functionality}
\logpage{[ 3, 11, 0 ]}
\hyperdef{L}{X78883D5E83B4425F}{}
{
  There is an exception to the rule that objects can only be modified if a
thread has write access to a region. A limited sets of objects can be modified
using the "bind once" family of functions. These allow the modifications of
objects to which a thread has read access in a limited fashion. 

 For reasons of implementation symmetry, these functions can also be used on
the atomic versions of these objects. 

 Implementation note: The functionality is not currently available for
component objects. 

\subsection{\textcolor{Chapter }{BindOnce}}
\logpage{[ 3, 11, 1 ]}\nobreak
\hyperdef{L}{X83AD36A68503CF70}{}
{\noindent\textcolor{FuncColor}{$\triangleright$\enspace\texttt{BindOnce({\mdseries\slshape obj, index, value})\index{BindOnce@\texttt{BindOnce}}
\label{BindOnce}
}\hfill{\scriptsize (function)}}\\


 \texttt{BindOnce} modifies \mbox{\texttt{\mdseries\slshape obj}}, which can be a positional object, atomic positional object, component
object, or atomic component object. It inspects \texttt{obj![index]} for the positional versions or \texttt{obj!.(index)} for the component versions. If the respective element is not yet bound, \mbox{\texttt{\mdseries\slshape value}} is assigned to that element. Otherwise, no modification happens. The test and
modification occur as one atomic step. The function returns the value of the
element; i.e. the old value if the element was bound and \mbox{\texttt{\mdseries\slshape value}} if it was unbound. 

 The intent of this function is to allow concurrent initialization of objects,
where multiple threads may attempt to set a value concurrently. Only one will
succeed; all threads can then use the return value of \texttt{BindOnce} as the definitive value of the element. It also allows for the lazy
initialization of objects in the read\texttt{\symbol{45}}only region. 

 The current thread needs to have at least read access to \mbox{\texttt{\mdseries\slshape obj}}, but does not require write access. }

 

\subsection{\textcolor{Chapter }{TestBindOnce}}
\logpage{[ 3, 11, 2 ]}\nobreak
\hyperdef{L}{X8673539F7DA79110}{}
{\noindent\textcolor{FuncColor}{$\triangleright$\enspace\texttt{TestBindOnce({\mdseries\slshape obj, index, value})\index{TestBindOnce@\texttt{TestBindOnce}}
\label{TestBindOnce}
}\hfill{\scriptsize (function)}}\\


 \texttt{TestBindOnce} works like \texttt{BindOnce} (\ref{BindOnce}), except that it returns \texttt{true} if the value could be bound and \texttt{false} otherwise. }

 

\subsection{\textcolor{Chapter }{BindOnceExpr}}
\logpage{[ 3, 11, 3 ]}\nobreak
\hyperdef{L}{X7BAEC41C87E1DC43}{}
{\noindent\textcolor{FuncColor}{$\triangleright$\enspace\texttt{BindOnceExpr({\mdseries\slshape obj, index, expr})\index{BindOnceExpr@\texttt{BindOnceExpr}}
\label{BindOnceExpr}
}\hfill{\scriptsize (function)}}\\


 \texttt{BindOnceExpr} works like \texttt{BindOnce} (\ref{BindOnce}), except that it evaluates the parameterless function \mbox{\texttt{\mdseries\slshape expr}} to determine the value. It will only evaluate \mbox{\texttt{\mdseries\slshape expr}} if the element is not bound. 

 For positional objects, the implementation works as follows: 
\begin{Verbatim}[commandchars=@|A,fontsize=\small,frame=single,label=Example]
  BindOnceExprPosObj := function(obj, index, expr)
    if not IsBound(obj![index]) then
      return BindOnce(obj, index, expr());
    else
      return obj![index]);
    fi;
  end;
\end{Verbatim}
 The implementation for component objects works analogously. 

 The intent is to avoid unnecessary computations if the value is already bound.
Note that this cannot be avoided entirely, because \texttt{obj![index]} or \texttt{obj!.(index)} can be bound while \mbox{\texttt{\mdseries\slshape expr}} is evaluated, but it can minimize such occurrences. }

 

\subsection{\textcolor{Chapter }{TestBindOnceExpr}}
\logpage{[ 3, 11, 4 ]}\nobreak
\hyperdef{L}{X81B15A9C8795DF59}{}
{\noindent\textcolor{FuncColor}{$\triangleright$\enspace\texttt{TestBindOnceExpr({\mdseries\slshape obj, index, expr})\index{TestBindOnceExpr@\texttt{TestBindOnceExpr}}
\label{TestBindOnceExpr}
}\hfill{\scriptsize (function)}}\\


 \texttt{TestBindOnceExpr} works like \texttt{BindOnceExpr} (\ref{BindOnceExpr}), except that it returns \texttt{true} if the value could be bound and \texttt{false} otherwise. }

 

\subsection{\textcolor{Chapter }{StrictBindOnce}}
\logpage{[ 3, 11, 5 ]}\nobreak
\hyperdef{L}{X7897092C86AE17D7}{}
{\noindent\textcolor{FuncColor}{$\triangleright$\enspace\texttt{StrictBindOnce({\mdseries\slshape obj, index, expr})\index{StrictBindOnce@\texttt{StrictBindOnce}}
\label{StrictBindOnce}
}\hfill{\scriptsize (function)}}\\


 \texttt{StrictBindOnce} works like \texttt{BindOnce} (\ref{BindOnce}), except that it raises an error if the element is already bound. This is
intended for cases where a read\texttt{\symbol{45}}only object is initialized,
but where another thread trying to initialize it concurrently would be an
error. }

 }

 }

 
\chapter{\textcolor{Chapter }{Console User Interface}}\label{Console User Interface}
\logpage{[ 4, 0, 0 ]}
\hyperdef{L}{X7AFF436381C319CD}{}
{
  HPC\texttt{\symbol{45}}GAP has a multi\texttt{\symbol{45}}threaded user
interface to assist with the development and debugging of concurrent programs.
This user interface is enabled by default; to disable it, and use the
single\texttt{\symbol{45}}threaded interface, GAP has to be started with the \texttt{\texttt{\symbol{45}}S} option. 
\section{\textcolor{Chapter }{Console UI commands}}\label{Console UI commands}
\logpage{[ 4, 1, 0 ]}
\hyperdef{L}{X7B958D607F9E0EF3}{}
{
  The console user interface provides the user with the option to control
threads by commands prefixed with an exclamation mark ("!"). Those commands
are listed below. 

 For ease of use, users only need to type as many letters of each commands so
that it can be unambiguously selected. Thus, the shell will recognize \texttt{!l} as an abbreviation for \texttt{!list}. 
\subsection{\textcolor{Chapter }{!shell [name]}}\label{shell}
\logpage{[ 4, 1, 1 ]}
\hyperdef{L}{X8149CF507A188BFB}{}
{
  Starts a new shell thread and switches to it. Optionally, a name for the
thread can be provided. 
\begin{Verbatim}[commandchars=@|A,fontsize=\small,frame=single,label=Example]
  @gapprompt|gap>A @gapinput|!shellA
  --- Switching to thread 4
  [4] gap>
\end{Verbatim}
 }

 
\subsection{\textcolor{Chapter }{!fork [name]}}\label{fork}
\logpage{[ 4, 1, 2 ]}
\hyperdef{L}{X8589D19484DBBA7B}{}
{
  Starts a new background shell thread. Optionally, a name for the thread can be
provided. 
\begin{Verbatim}[commandchars=@|A,fontsize=\small,frame=single,label=Example]
  @gapprompt|gap>A @gapinput|!forkA
  --- Created new thread 5
\end{Verbatim}
 }

 
\subsection{\textcolor{Chapter }{!list}}\label{!list}
\logpage{[ 4, 1, 3 ]}
\hyperdef{L}{X7E57FE757CBAD4BC}{}
{
  List all current threads that are interacting with the user. This does not
list threads created with \texttt{CreateThread()} that have not entered a break loop. 
\begin{Verbatim}[commandchars=@|A,fontsize=\small,frame=single,label=Example]
  @gapprompt|gap>A @gapinput|!listA
  --- Thread 0 [0]
  --- Thread 4 [4]
  --- Thread 5 [5] (pending output)
\end{Verbatim}
 }

 
\subsection{\textcolor{Chapter }{!kill id}}\label{kill}
\logpage{[ 4, 1, 4 ]}
\hyperdef{L}{X7C2721208570CF5D}{}
{
 Terminates the specified thread. }

 
\subsection{\textcolor{Chapter }{!break id}}\label{break}
\logpage{[ 4, 1, 5 ]}
\hyperdef{L}{X86052CA284DD8E8E}{}
{
 Makes the specified thread enter a break loop. }

 
\subsection{\textcolor{Chapter }{!name [id] name}}\label{name}
\logpage{[ 4, 1, 6 ]}
\hyperdef{L}{X8022D93682DD8ABE}{}
{
 Give the thread with the numerical identifier or name \texttt{id} the name \texttt{name}. 
\begin{Verbatim}[commandchars=@|A,fontsize=\small,frame=single,label=Example]
  @gapprompt|gap>A @gapinput|!name 5 testA
  @gapprompt|gap>A @gapinput|!listA
  --- Thread 0 [0]
  --- Thread 4 [4]
  --- Thread test [5] (pending output)
\end{Verbatim}
 }

 
\subsection{\textcolor{Chapter }{!info id}}\label{info}
\logpage{[ 4, 1, 7 ]}
\hyperdef{L}{X810ED84D848681B6}{}
{
  Provide information about the thread with the numerical identifier or name \texttt{id}. \emph{Not yet implemented}. }

 
\subsection{\textcolor{Chapter }{!hide [id|*]}}\label{hide}
\logpage{[ 4, 1, 8 ]}
\hyperdef{L}{X847F80AB873516DC}{}
{
  Hide output from the thread with the numerical identifier or name \texttt{id} when it is not the foreground thread. If no thread is specified, make this the
default behavior for future threads. }

 
\subsection{\textcolor{Chapter }{!watch [id|*]}}\label{watch}
\logpage{[ 4, 1, 9 ]}
\hyperdef{L}{X87515BA77E6327EE}{}
{
  Show output from the thread with the numerical identifier or name \texttt{id} even when it is not the foreground thread. If no thread is specified, make
this the default behavior for future threads. }

 
\subsection{\textcolor{Chapter }{!keep num}}\label{keep}
\logpage{[ 4, 1, 10 ]}
\hyperdef{L}{X8295E57782EEF3ED}{}
{
  Keep \texttt{num} lines of output from each thread. }

 
\subsection{\textcolor{Chapter }{!prompt (id|*) string}}\label{prompt}
\logpage{[ 4, 1, 11 ]}
\hyperdef{L}{X8305DFEA87C30627}{}
{
  Set the prompt for the specified thread (or for all newly created threads if \texttt{*} was specified) to be \texttt{string}. If the string contains the pattern \texttt{id}, it is replaced with the numerical id of the thread; if it contains the
pattern \texttt{name}, it is replaced with the name of the thread; if the thread has no name, the
numerical id is displayed instead. }

 
\subsection{\textcolor{Chapter }{!prefix (id|*) string}}\label{prefix}
\logpage{[ 4, 1, 12 ]}
\hyperdef{L}{X8484498A787248A2}{}
{
  Prefix the output from the specified thread (or for all newly created threads
if \texttt{*} was specified) with \texttt{string}. The same substitution rules as for the \texttt{!prompt} command apply. }

 
\subsection{\textcolor{Chapter }{!select id}}\label{select}
\logpage{[ 4, 1, 13 ]}
\hyperdef{L}{X80CE7A5B874A8D6B}{}
{
  Make the specified thread the foreground thread. 
\begin{Verbatim}[commandchars=@|A,fontsize=\small,frame=single,label=Example]
  @gapprompt|gap>A @gapinput|!select 4A
  @gapprompt|gap>A @gapinput|!select 4A
  --- Switching to thread 4
  [4] gap>
\end{Verbatim}
 }

 
\subsection{\textcolor{Chapter }{!next}}\label{next}
\logpage{[ 4, 1, 14 ]}
\hyperdef{L}{X7A5980F782780572}{}
{
  Make the next thread in numerical order the foreground thread. }

 
\subsection{\textcolor{Chapter }{!previous}}\label{previous}
\logpage{[ 4, 1, 15 ]}
\hyperdef{L}{X838676D07FBD27CE}{}
{
  Make the previous thread in numerical order the foreground thread. }

 
\subsection{\textcolor{Chapter }{!replay num [id]}}\label{replay}
\logpage{[ 4, 1, 16 ]}
\hyperdef{L}{X7B6441957CFB191E}{}
{
  Display the last \texttt{num} lines of output of the specified thread. If no thread was specified, display
the last \texttt{num} lines of the current foreground thread. }

 
\subsection{\textcolor{Chapter }{!id}}\label{id}
\logpage{[ 4, 1, 17 ]}
\hyperdef{L}{X850D86EE8203F94D}{}
{
  \texttt{!id} is a shortcut for \texttt{!select id}. }

 
\subsection{\textcolor{Chapter }{!source file}}\label{source}
\logpage{[ 4, 1, 18 ]}
\hyperdef{L}{X84FE74AA7C34C493}{}
{
  Read commands from file \texttt{file}. }

 
\subsection{\textcolor{Chapter }{!alias shortcut expansion}}\label{alias}
\logpage{[ 4, 1, 19 ]}
\hyperdef{L}{X84BA0EF7803E2920}{}
{
  Create an alias. After defining the alias, \texttt{!shortcut 'rest of line'} will be replaced with \texttt{!expansion 'rest of line'}. }

 
\subsection{\textcolor{Chapter }{!unalias shortcut}}\label{unalias}
\logpage{[ 4, 1, 20 ]}
\hyperdef{L}{X80E7398679C4B87C}{}
{
  Removes the specified alias. }

 
\subsection{\textcolor{Chapter }{!eval expr}}\label{eval}
\logpage{[ 4, 1, 21 ]}
\hyperdef{L}{X79D5CE9A86D5FC6A}{}
{
 Evaluates \texttt{expr} as a command. }

 
\subsection{\textcolor{Chapter }{!run function string}}\label{run}
\logpage{[ 4, 1, 22 ]}
\hyperdef{L}{X831807A67C6B4B2D}{}
{
  Calls the function with name \texttt{function}, passing it the single argument \texttt{string} as a GAP string. }

 }

 
\section{\textcolor{Chapter }{GAP functions to access the Shell UI}}\label{GAP functions to access the Shell UI}
\logpage{[ 4, 2, 0 ]}
\hyperdef{L}{X876007F6862BA1BD}{}
{
  There are several functions to access the basic functionality of the shell
user interface. Other than \texttt{TextUIRegisterCommand} (\ref{TextUIRegisterCommand}), they can only be called from within a registered command. 

 Threads can be specified either by their numerical identifier or by their name
(as a string). The empty string can be used to specify the current foreground
thread. 

\subsection{\textcolor{Chapter }{TextUIRegisterCommand}}
\logpage{[ 4, 2, 1 ]}\nobreak
\hyperdef{L}{X84154E44780A3402}{}
{\noindent\textcolor{FuncColor}{$\triangleright$\enspace\texttt{TextUIRegisterCommand({\mdseries\slshape name, func})\index{TextUIRegisterCommand@\texttt{TextUIRegisterCommand}}
\label{TextUIRegisterCommand}
}\hfill{\scriptsize (function)}}\\


 Registers the command \texttt{!name} with the shell UI. It will call {\textless}func{\textgreater} with the rest of
the command line passed as a string argument when typed. }

 

\subsection{\textcolor{Chapter }{TextUIForegroundThread}}
\logpage{[ 4, 2, 2 ]}\nobreak
\hyperdef{L}{X85312D647966343E}{}
{\noindent\textcolor{FuncColor}{$\triangleright$\enspace\texttt{TextUIForegroundThread({\mdseries\slshape })\index{TextUIForegroundThread@\texttt{TextUIForegroundThread}}
\label{TextUIForegroundThread}
}\hfill{\scriptsize (function)}}\\


 Returns the numerical identifier of the current foreground thread. }

 

\subsection{\textcolor{Chapter }{TextUIForegroundThreadName}}
\logpage{[ 4, 2, 3 ]}\nobreak
\hyperdef{L}{X7AD377F882A1EEC8}{}
{\noindent\textcolor{FuncColor}{$\triangleright$\enspace\texttt{TextUIForegroundThreadName({\mdseries\slshape })\index{TextUIForegroundThreadName@\texttt{TextUIForegroundThreadName}}
\label{TextUIForegroundThreadName}
}\hfill{\scriptsize (function)}}\\


 Returns the name of the current foreground thread or \texttt{fail} if the current foreground thread has no name. }

 

\subsection{\textcolor{Chapter }{TextUISelectThread}}
\logpage{[ 4, 2, 4 ]}\nobreak
\hyperdef{L}{X8117AC7980300EF5}{}
{\noindent\textcolor{FuncColor}{$\triangleright$\enspace\texttt{TextUISelectThread({\mdseries\slshape id})\index{TextUISelectThread@\texttt{TextUISelectThread}}
\label{TextUISelectThread}
}\hfill{\scriptsize (function)}}\\


 Makes \mbox{\texttt{\mdseries\slshape id}} the current foreground thread. Returns \texttt{true} or \texttt{false} to indicate success. }

 

\subsection{\textcolor{Chapter }{TextUIOutputHistory}}
\logpage{[ 4, 2, 5 ]}\nobreak
\hyperdef{L}{X7C9543F07F0FECF5}{}
{\noindent\textcolor{FuncColor}{$\triangleright$\enspace\texttt{TextUIOutputHistory({\mdseries\slshape id, count})\index{TextUIOutputHistory@\texttt{TextUIOutputHistory}}
\label{TextUIOutputHistory}
}\hfill{\scriptsize (function)}}\\


 Returns the last \mbox{\texttt{\mdseries\slshape count}} lines of the thread specified by \mbox{\texttt{\mdseries\slshape id}} (which can be a numerical identifier or a name). Returns \texttt{fail} if there is no such thread. }

 

\subsection{\textcolor{Chapter }{TextUISetOutputHistoryLength}}
\logpage{[ 4, 2, 6 ]}\nobreak
\hyperdef{L}{X839D85B97B23BC3C}{}
{\noindent\textcolor{FuncColor}{$\triangleright$\enspace\texttt{TextUISetOutputHistoryLength({\mdseries\slshape length})\index{TextUISetOutputHistoryLength@\texttt{TextUISetOutputHistoryLength}}
\label{TextUISetOutputHistoryLength}
}\hfill{\scriptsize (function)}}\\


 By default, retain \mbox{\texttt{\mdseries\slshape length}} lines of output history from each thread. }

 

\subsection{\textcolor{Chapter }{TextUINewSession}}
\logpage{[ 4, 2, 7 ]}\nobreak
\hyperdef{L}{X81AB547681B0A2C8}{}
{\noindent\textcolor{FuncColor}{$\triangleright$\enspace\texttt{TextUINewSession({\mdseries\slshape foreground, name})\index{TextUINewSession@\texttt{TextUINewSession}}
\label{TextUINewSession}
}\hfill{\scriptsize (function)}}\\


 Creates a new shell thread. Here, \mbox{\texttt{\mdseries\slshape foreground}} is a boolean variable specifying whether it should be made the new foreground
thread and \mbox{\texttt{\mdseries\slshape name}} is the name of the thread. The empty string can be used to leave the thread
without a name. }

 

\subsection{\textcolor{Chapter }{TextUIRunCommand}}
\logpage{[ 4, 2, 8 ]}\nobreak
\hyperdef{L}{X7B69D7177D7742D2}{}
{\noindent\textcolor{FuncColor}{$\triangleright$\enspace\texttt{TextUIRunCommand({\mdseries\slshape command})\index{TextUIRunCommand@\texttt{TextUIRunCommand}}
\label{TextUIRunCommand}
}\hfill{\scriptsize (function)}}\\


 Run the command denoted by \mbox{\texttt{\mdseries\slshape command}} as though a user had typed it. The command must not contain a newline
character. }

 

\subsection{\textcolor{Chapter }{TextUIWritePrompt}}
\logpage{[ 4, 2, 9 ]}\nobreak
\hyperdef{L}{X8532502D78CE5475}{}
{\noindent\textcolor{FuncColor}{$\triangleright$\enspace\texttt{TextUIWritePrompt({\mdseries\slshape })\index{TextUIWritePrompt@\texttt{TextUIWritePrompt}}
\label{TextUIWritePrompt}
}\hfill{\scriptsize (function)}}\\


 Display a prompt for the current thread. }

 }

 }

 
\chapter{\textcolor{Chapter }{Atomic objects}}\label{Atomic objects}
\logpage{[ 5, 0, 0 ]}
\hyperdef{L}{X7A5E986A790A5A18}{}
{
  HPC\texttt{\symbol{45}}GAP provides a number of atomic object types. These can
be accessed by multiple threads concurrently without requiring explicit
synchronization, but can have non\texttt{\symbol{45}}deterministic behavior
for complex operations. Atomic lists are fixed\texttt{\symbol{45}}size lists;
they can be assigned to and read from like normal plain lists. Atomic records
are atomic versions of plain records. Unlike plain records, though, it is not
possible to delete elements from an atomic record. The primary use of atomic
lists and records is to facilitate storing the result of idempotent operations
and to support certain low\texttt{\symbol{45}}level operations. Atomic lists
and records can have three different replacement policies:
write\texttt{\symbol{45}}once, strict write\texttt{\symbol{45}}once, and
rewritable. The replacement policy determines whether an already assigned
element can be changed. The write\texttt{\symbol{45}}once policy allows
elements to be assigned only once, with subsequent assignments being ignored;
the strict write\texttt{\symbol{45}}once policy allows elements also to be
assigned only once, but subsequent assignments will raise an error; the
rewritable policy allows elements to be assigned different values repeatedly.
The default for new atomic objects is to be rewritable.
Thread\texttt{\symbol{45}}local records are variants of plain records that are
replicated on a per\texttt{\symbol{45}}thread basis. 
\section{\textcolor{Chapter }{Atomic lists}}\label{Atomic lists}
\logpage{[ 5, 1, 0 ]}
\hyperdef{L}{X79DFD10A80868799}{}
{
  Atomic lists are created using the \texttt{AtomicList} or \texttt{FixedAtomicList} functions. After creation, they can be used exactly like any other list,
except that atomic lists created with \texttt{FixedAtomicList} cannot be resized. Their contents can also be read as normal plain lists using \texttt{FromAtomicList}. 
\begin{Verbatim}[commandchars=!@|,fontsize=\small,frame=single,label=Example]
  !gapprompt@gap>| !gapinput@a := AtomicList([1,2,4]);|
  <atomic list of size 3>
  !gapprompt@gap>| !gapinput@WaitTask(RunTask(function() a[1] := a[1] + a[2]; end));|
  !gapprompt@gap>| !gapinput@a[1];|
  3
  !gapprompt@gap>| !gapinput@FromAtomicList(a);|
  [ 3, 2, 4 ]
\end{Verbatim}
 Because multiple threads can read and write the list concurrently without
synchronization, the results of modifying the list may be
non\texttt{\symbol{45}}deterministic. It is faster to write to fixed atomic
lists than to a resizable atomic list. 

\subsection{\textcolor{Chapter }{AtomicList}}
\logpage{[ 5, 1, 1 ]}\nobreak
\hyperdef{L}{X7F1353B58414D4C3}{}
{\noindent\textcolor{FuncColor}{$\triangleright$\enspace\texttt{AtomicList({\mdseries\slshape list})\index{AtomicList@\texttt{AtomicList}}
\label{AtomicList}
}\hfill{\scriptsize (function)}}\\
\noindent\textcolor{FuncColor}{$\triangleright$\enspace\texttt{AtomicList({\mdseries\slshape count, obj})\index{AtomicList@\texttt{AtomicList}!for a count and an object}
\label{AtomicList:for a count and an object}
}\hfill{\scriptsize (function)}}\\


 \texttt{AtomicList} is used to create a new atomic list. It takes either a plain list as an
argument, in which case it will create a new atomic list of the same size,
populated by the same elements; or it takes a count and an object argument. In
that case, it creates an atomic list with \mbox{\texttt{\mdseries\slshape count}} elements, each set to the value of \mbox{\texttt{\mdseries\slshape obj}}. 
\begin{Verbatim}[commandchars=!@|,fontsize=\small,frame=single,label=Example]
  !gapprompt@gap>| !gapinput@al := AtomicList([3, 1, 4]);|
  <atomic list of size 3>
  !gapprompt@gap>| !gapinput@al[3];|
  4
  !gapprompt@gap>| !gapinput@al := AtomicList(10, `"alpha");|
  <atomic list of size 10>
  !gapprompt@gap>| !gapinput@al[3];|
  "alpha"
  !gapprompt@gap>| !gapinput@WaitTask(RunTask(function() al[3] := `"beta"; end));|
  !gapprompt@gap>| !gapinput@al[3];|
  "beta"
\end{Verbatim}
 }

 

\subsection{\textcolor{Chapter }{FixedAtomicList}}
\logpage{[ 5, 1, 2 ]}\nobreak
\hyperdef{L}{X8279EC4483E0562C}{}
{\noindent\textcolor{FuncColor}{$\triangleright$\enspace\texttt{FixedAtomicList({\mdseries\slshape list})\index{FixedAtomicList@\texttt{FixedAtomicList}}
\label{FixedAtomicList}
}\hfill{\scriptsize (function)}}\\
\noindent\textcolor{FuncColor}{$\triangleright$\enspace\texttt{FixedAtomicList({\mdseries\slshape count, obj})\index{FixedAtomicList@\texttt{FixedAtomicList}!for a count and an object}
\label{FixedAtomicList:for a count and an object}
}\hfill{\scriptsize (function)}}\\


 \texttt{FixedAtomicList} works like \texttt{AtomicList} (\ref{AtomicList}) except that the resulting list cannot be resized. }

 

\subsection{\textcolor{Chapter }{MakeFixedAtomicList}}
\logpage{[ 5, 1, 3 ]}\nobreak
\hyperdef{L}{X807687A579422D3B}{}
{\noindent\textcolor{FuncColor}{$\triangleright$\enspace\texttt{MakeFixedAtomicList({\mdseries\slshape list})\index{MakeFixedAtomicList@\texttt{MakeFixedAtomicList}}
\label{MakeFixedAtomicList}
}\hfill{\scriptsize (function)}}\\


 \texttt{MakeFixedAtomicList} turns a resizable atomic list into a fixed atomic list. 
\begin{Verbatim}[commandchars=!@|,fontsize=\small,frame=single,label=Example]
  !gapprompt@gap>| !gapinput@a := AtomicList([99]);|
  <atomic list of size 1>
  !gapprompt@gap>| !gapinput@a[2] := 100;|
  100
  !gapprompt@gap>| !gapinput@MakeFixedAtomicList(a);|
  <fixed atomic list of size 2>
  !gapprompt@gap>| !gapinput@a[3] := 101;|
  Error, Atomic List Element: <pos>=3 is an invalid index for <list>
\end{Verbatim}
 }

 

\subsection{\textcolor{Chapter }{FromAtomicList}}
\logpage{[ 5, 1, 4 ]}\nobreak
\hyperdef{L}{X81AF8061871DA27C}{}
{\noindent\textcolor{FuncColor}{$\triangleright$\enspace\texttt{FromAtomicList({\mdseries\slshape atomic{\textunderscore}list})\index{FromAtomicList@\texttt{FromAtomicList}}
\label{FromAtomicList}
}\hfill{\scriptsize (function)}}\\


 \texttt{FromAtomicList} returns a plain list containing the same elements as \mbox{\texttt{\mdseries\slshape atomic{\textunderscore}list}} at the time of the call. Because other threads can write concurrently to that
list, the result is not guaranteed to be deterministic. 
\begin{Verbatim}[commandchars=!@|,fontsize=\small,frame=single,label=Example]
  !gapprompt@gap>| !gapinput@al := AtomicList([10, 20, 30]);;|
  !gapprompt@gap>| !gapinput@WaitTask(RunTask(function() al[2] := 40; end));|
  !gapprompt@gap>| !gapinput@FromAtomicList(al);|
  [ 10, 40, 30 ]
\end{Verbatim}
 }

 

\subsection{\textcolor{Chapter }{ATOMIC{\textunderscore}ADDITION}}
\logpage{[ 5, 1, 5 ]}\nobreak
\hyperdef{L}{X795EE57983074ADF}{}
{\noindent\textcolor{FuncColor}{$\triangleright$\enspace\texttt{ATOMIC{\textunderscore}ADDITION({\mdseries\slshape atomic{\textunderscore}list, index, value})\index{ATOMIC{\textunderscore}ADDITION@\texttt{ATOMIC{\textunderscore}ADDITION}}
\label{ATOMICuScoreADDITION}
}\hfill{\scriptsize (function)}}\\


 \texttt{ATOMIC{\textunderscore}ADDITION} is a low\texttt{\symbol{45}}level operation that atomically adds \mbox{\texttt{\mdseries\slshape value}} to \mbox{\texttt{\mdseries\slshape atomic{\textunderscore}list[index]}}. It returns the value of \mbox{\texttt{\mdseries\slshape atomic{\textunderscore}list[index]}} after the addition has been performed. 
\begin{Verbatim}[commandchars=!@|,fontsize=\small,frame=single,label=Example]
  !gapprompt@gap>| !gapinput@al := FixedAtomicList([4,5,6]);;|
  !gapprompt@gap>| !gapinput@ATOMIC_ADDITION(al, 2, 7);|
  12
  !gapprompt@gap>| !gapinput@FromAtomicList(al);|
  [ 4, 12, 6 ]
\end{Verbatim}
 }

 

\subsection{\textcolor{Chapter }{COMPARE{\textunderscore}AND{\textunderscore}SWAP}}
\logpage{[ 5, 1, 6 ]}\nobreak
\hyperdef{L}{X7847EAA284A02AEE}{}
{\noindent\textcolor{FuncColor}{$\triangleright$\enspace\texttt{COMPARE{\textunderscore}AND{\textunderscore}SWAP({\mdseries\slshape atomic{\textunderscore}list, index, old, new})\index{COMPARE{\textunderscore}AND{\textunderscore}SWAP@\texttt{COM}\-\texttt{P}\-\texttt{A}\-\texttt{R}\-\texttt{E{\textunderscore}}\-\texttt{A}\-\texttt{N}\-\texttt{D{\textunderscore}}\-\texttt{SWAP}}
\label{COMPAREuScoreANDuScoreSWAP}
}\hfill{\scriptsize (function)}}\\


 \texttt{COMPARE{\textunderscore}AND{\textunderscore}SWAP} is an atomic operation. It atomically compares \mbox{\texttt{\mdseries\slshape atomic{\textunderscore}list[index]}} to \mbox{\texttt{\mdseries\slshape old}} and, if they are identical, replaces the value (in the same atomic step) with \mbox{\texttt{\mdseries\slshape new}}. It returns true if the replacement took place, false otherwise. 

 The primary use of \texttt{COMPARE{\textunderscore}AND{\textunderscore}SWAP} is to implement certain concurrency primitives; most programmers will not need
to use it. }

 }

 
\section{\textcolor{Chapter }{Atomic records and component objects}}\label{Atomic records and component objects}
\logpage{[ 5, 2, 0 ]}
\hyperdef{L}{X823767367D6CF4FD}{}
{
  Atomic records are atomic counterparts to plain records. They support
assignment to individual record fields, and conversion to and from plain
records. 

 Assignment semantics can be specified on a per\texttt{\symbol{45}}record basis
if the assigned record field is already populated, allowing either an
overwrite, keeping the existing value, or raising an error. 

 It is not possible to unbind atomic record elements. 

 Like plain records, atomic records can be converted to component objects using \texttt{Objectify}. 

\subsection{\textcolor{Chapter }{AtomicRecord}}
\logpage{[ 5, 2, 1 ]}\nobreak
\hyperdef{L}{X87EA343179BE4BFE}{}
{\noindent\textcolor{FuncColor}{$\triangleright$\enspace\texttt{AtomicRecord({\mdseries\slshape capacity})\index{AtomicRecord@\texttt{AtomicRecord}}
\label{AtomicRecord}
}\hfill{\scriptsize (function)}}\\
\noindent\textcolor{FuncColor}{$\triangleright$\enspace\texttt{AtomicRecord({\mdseries\slshape record})\index{AtomicRecord@\texttt{AtomicRecord}!for a record}
\label{AtomicRecord:for a record}
}\hfill{\scriptsize (function)}}\\


 \texttt{AtomicRecord} is used to create a new atomic record. Its single optional argument is either
a positive integer, denoting the intended capacity (i.e., number of elements
to be held) of the record, in which case a new empty atomic record with that
initial capacity will be created. Alternatively, the caller can provide a
plain record with which to initially populate the atomic record. 
\begin{Verbatim}[commandchars=!@|,fontsize=\small,frame=single,label=Example]
  !gapprompt@gap>| !gapinput@r := AtomicRecord(rec( x := 2 ));|
  <atomic record 1/2 full>
  !gapprompt@gap>| !gapinput@r.y := 3;|
  3
  !gapprompt@gap>| !gapinput@TaskResult(RunTask(function() return r.x + r.y; end));|
  5
  !gapprompt@gap>| !gapinput@[ r.x, r.y ];|
  [ 2, 3 ]
\end{Verbatim}
 Any atomic record can later grow beyond its initial capacity. There is no
limit to the number of elements it can hold other than available memory. }

 

\subsection{\textcolor{Chapter }{FromAtomicRecord}}
\logpage{[ 5, 2, 2 ]}\nobreak
\hyperdef{L}{X81E6B21C7A1AC7FE}{}
{\noindent\textcolor{FuncColor}{$\triangleright$\enspace\texttt{FromAtomicRecord({\mdseries\slshape record})\index{FromAtomicRecord@\texttt{FromAtomicRecord}}
\label{FromAtomicRecord}
}\hfill{\scriptsize (function)}}\\


 \texttt{FromAtomicRecord} returns a plain record copy of the atomic record \mbox{\texttt{\mdseries\slshape record}}. This copy is shallow; elements of \mbox{\texttt{\mdseries\slshape record}} will not also be copied. 
\begin{Verbatim}[commandchars=!@|,fontsize=\small,frame=single,label=Example]
  !gapprompt@gap>| !gapinput@r := AtomicRecord();;|
  !gapprompt@gap>| !gapinput@r.x := 1;; r.y := 2;; r.z := 3;;|
  !gapprompt@gap>| !gapinput@FromAtomicRecord(r);|
  rec( x := 1, y := 2, z := 3 )
\end{Verbatim}
 }

 }

 
\section{\textcolor{Chapter }{Replacement policy functions}}\label{Replacement policy functions}
\logpage{[ 5, 3, 0 ]}
\hyperdef{L}{X78EDB1DB79473F53}{}
{
  There are three functions that set the replacement policy of an atomic object.
All three can also be used with plain lists and records, in which case an
atomic version of the list or record is first created. This allows programmers
to elide \texttt{AtomicList} (\ref{AtomicList}) and \texttt{AtomicRecord} (\ref{AtomicRecord}) calls when the next step is to change their policy. 

\subsection{\textcolor{Chapter }{MakeWriteOnceAtomic}}
\logpage{[ 5, 3, 1 ]}\nobreak
\hyperdef{L}{X7B76F2AD818492BC}{}
{\noindent\textcolor{FuncColor}{$\triangleright$\enspace\texttt{MakeWriteOnceAtomic({\mdseries\slshape obj})\index{MakeWriteOnceAtomic@\texttt{MakeWriteOnceAtomic}}
\label{MakeWriteOnceAtomic}
}\hfill{\scriptsize (function)}}\\


 \texttt{MakeWriteOnceAtomic} takes a list, record, atomic list, atomic record, atomic positional object, or
atomic component object as its argument. If the argument is a
non\texttt{\symbol{45}}atomic list or record, then the function first creates
an atomic copy of the argument. The function then changes the replacement
policy of the object to write\texttt{\symbol{45}}once: if an element of the
object is already bound, then further attempts to assign to it will be
ignored. }

 

\subsection{\textcolor{Chapter }{MakeStrictWriteOnceAtomic}}
\logpage{[ 5, 3, 2 ]}\nobreak
\hyperdef{L}{X806C3F0681B847A0}{}
{\noindent\textcolor{FuncColor}{$\triangleright$\enspace\texttt{MakeStrictWriteOnceAtomic({\mdseries\slshape obj})\index{MakeStrictWriteOnceAtomic@\texttt{MakeStrictWriteOnceAtomic}}
\label{MakeStrictWriteOnceAtomic}
}\hfill{\scriptsize (function)}}\\


 \texttt{MakeStrictWriteOnceAtomic} works like \texttt{MakeWriteOnceAtomic} (\ref{MakeWriteOnceAtomic}), except that the replacement policy is being changed to being strict
write\texttt{\symbol{45}}once: if an element is already bound, then further
attempts to assign to it will raise an error. }

 

\subsection{\textcolor{Chapter }{MakeReadWriteAtomic}}
\logpage{[ 5, 3, 3 ]}\nobreak
\hyperdef{L}{X7E7591EC780F3988}{}
{\noindent\textcolor{FuncColor}{$\triangleright$\enspace\texttt{MakeReadWriteAtomic({\mdseries\slshape obj})\index{MakeReadWriteAtomic@\texttt{MakeReadWriteAtomic}}
\label{MakeReadWriteAtomic}
}\hfill{\scriptsize (function)}}\\


 \texttt{MakeReadWriteAtomic} is the inverse of \texttt{MakeWriteOnceAtomic} (\ref{MakeWriteOnceAtomic}) and \texttt{MakeStrictWriteOnceAtomic} (\ref{MakeStrictWriteOnceAtomic}) in that the replacement policy is being changed to being rewritable: Elements
can be replaced even if they are already bound. }

 }

 
\section{\textcolor{Chapter }{Thread\texttt{\symbol{45}}local records}}\label{Thread-local records}
\logpage{[ 5, 4, 0 ]}
\hyperdef{L}{X7DCE3B4D7D26407F}{}
{
  Thread\texttt{\symbol{45}}local records allow an easy way to have a separate
copy of a record for each individual thread that is accessed by the same name
in each thread. 
\begin{Verbatim}[commandchars=!@|,fontsize=\small,frame=single,label=Example]
  !gapprompt@gap>| !gapinput@r := ThreadLocalRecord();; # create new thread-local record|
  !gapprompt@gap>| !gapinput@r.x := 99;;|
  !gapprompt@gap>| !gapinput@WaitThread( CreateThread( function()|
  !gapprompt@>| !gapinput@                             r.x := 100;|
  !gapprompt@>| !gapinput@                             Display(r.x);|
  !gapprompt@>| !gapinput@                             end ) );|
  100
  !gapprompt@gap>| !gapinput@r.x;|
  99
\end{Verbatim}
 As can be seen above, even though \texttt{r.x} is overwritten in the second thread, it does not affect the value of \texttt{r.x| in the first thread} 

\subsection{\textcolor{Chapter }{ThreadLocalRecord}}
\logpage{[ 5, 4, 1 ]}\nobreak
\hyperdef{L}{X7BE7D67586FC91E6}{}
{\noindent\textcolor{FuncColor}{$\triangleright$\enspace\texttt{ThreadLocalRecord({\mdseries\slshape [defaults[, constructors]]})\index{ThreadLocalRecord@\texttt{ThreadLocalRecord}}
\label{ThreadLocalRecord}
}\hfill{\scriptsize (function)}}\\


 \texttt{ThreadLocalRecord} creates a new thread\texttt{\symbol{45}}local record. It accepts up to two
initial arguments. The \mbox{\texttt{\mdseries\slshape defaults}} argument is a record of default values with which each
thread\texttt{\symbol{45}}local copy is initially populated (this happens on
demand, so values are not actually read until needed). The second argument is
a record of constructors; parameterless functions that return an initial value
for the respective element. Constructors are evaluated only once per thread
and only if the respective element is accessed without having previously been
assigned a value. 
\begin{Verbatim}[commandchars=!@|,fontsize=\small,frame=single,label=Example]
  !gapprompt@gap>| !gapinput@r := ThreadLocalRecord( rec(x := 99),|
  !gapprompt@>| !gapinput@     rec(y := function() return 101; end));;|
  !gapprompt@gap>| !gapinput@r.x;|
  99
  !gapprompt@gap>| !gapinput@r.y;|
  101
  !gapprompt@gap>| !gapinput@TaskResult(RunTask(function() return r.x; end));|
  99
  !gapprompt@gap>| !gapinput@TaskResult(RunTask(function() return r.y; end));|
  101
\end{Verbatim}
 }

 

\subsection{\textcolor{Chapter }{SetTLDefault}}
\logpage{[ 5, 4, 2 ]}\nobreak
\hyperdef{L}{X788AF1CC8111798C}{}
{\noindent\textcolor{FuncColor}{$\triangleright$\enspace\texttt{SetTLDefault({\mdseries\slshape record, name, value})\index{SetTLDefault@\texttt{SetTLDefault}}
\label{SetTLDefault}
}\hfill{\scriptsize (function)}}\\


 \texttt{SetTLDefault} can be used to set the default value of a record field after its creation.
Here, \mbox{\texttt{\mdseries\slshape record}} is a thread\texttt{\symbol{45}}local record, \mbox{\texttt{\mdseries\slshape name}} is the string of the field name, and \mbox{\texttt{\mdseries\slshape value}} is the value. 
\begin{Verbatim}[commandchars=!@|,fontsize=\small,frame=single,label=Example]
  !gapprompt@gap>| !gapinput@r := ThreadLocalRecord();;|
  !gapprompt@gap>| !gapinput@SetTLDefault(r, "x", 314);|
  !gapprompt@gap>| !gapinput@r.x;|
  314
  !gapprompt@gap>| !gapinput@TaskResult(RunTask(function() return r.x; end));|
  314
\end{Verbatim}
 }

 

\subsection{\textcolor{Chapter }{SetTLConstructor}}
\logpage{[ 5, 4, 3 ]}\nobreak
\hyperdef{L}{X853A898384888B85}{}
{\noindent\textcolor{FuncColor}{$\triangleright$\enspace\texttt{SetTLConstructor({\mdseries\slshape record, name, func})\index{SetTLConstructor@\texttt{SetTLConstructor}}
\label{SetTLConstructor}
}\hfill{\scriptsize (function)}}\\


 \texttt{SetTLConstructor} can be used to set the constructor of a thread\texttt{\symbol{45}}local record
field after its creation, similar to \texttt{SetTLDefault} (\ref{SetTLDefault}). 
\begin{Verbatim}[commandchars=!@|,fontsize=\small,frame=single,label=Example]
  !gapprompt@gap>| !gapinput@r := ThreadLocalRecord();;|
  !gapprompt@gap>| !gapinput@SetTLConstructor(r, "x", function() return 2718; end);|
  !gapprompt@gap>| !gapinput@r.x;|
  2718
  !gapprompt@gap>| !gapinput@TaskResult(RunTask(function() return r.x; end));|
  2718
\end{Verbatim}
 }

 }

 }

 
\chapter{\textcolor{Chapter }{Thread functions}}\label{Thread functions}
\logpage{[ 6, 0, 0 ]}
\hyperdef{L}{X841085357B854E1C}{}
{
  HPC\texttt{\symbol{45}}GAP has low\texttt{\symbol{45}}level functionality to
support explicit creation of threads. In practice, programmers should use
higher\texttt{\symbol{45}}level functionality, such as tasks, to describe
concurrency. The thread functions described here exist to facilitate the
construction of higher level libraries and are not meant to be used directly. 
\section{\textcolor{Chapter }{Thread functions}}\label{section:Thread functions}
\logpage{[ 6, 1, 0 ]}
\hyperdef{L}{X841085357B854E1C}{}
{
  

\subsection{\textcolor{Chapter }{CreateThread}}
\logpage{[ 6, 1, 1 ]}\nobreak
\hyperdef{L}{X7A541B6C84980B7D}{}
{\noindent\textcolor{FuncColor}{$\triangleright$\enspace\texttt{CreateThread({\mdseries\slshape func[, arg1, ..., argn]})\index{CreateThread@\texttt{CreateThread}}
\label{CreateThread}
}\hfill{\scriptsize (function)}}\\


 New threads are created with the function \texttt{CreateThread}. The thread takes at least one function as its argument that it will call in
the newly created thread; it also accepts zero or more parameters that will be
passed to that function. 

 The function returns a thread object describing the thread. 

 Only a finite number of threads can be active at a time (that limit is
system\texttt{\symbol{45}}dependent). To reclaim the resources occupied by one
thread, use the \texttt{WaitThread} (\ref{WaitThread}) function. }

 

\subsection{\textcolor{Chapter }{WaitThread}}
\logpage{[ 6, 1, 2 ]}\nobreak
\hyperdef{L}{X81DFC84779CD1C1E}{}
{\noindent\textcolor{FuncColor}{$\triangleright$\enspace\texttt{WaitThread({\mdseries\slshape threadID})\index{WaitThread@\texttt{WaitThread}}
\label{WaitThread}
}\hfill{\scriptsize (function)}}\\


 The \texttt{WaitThread} function waits for the thread identified by \mbox{\texttt{\mdseries\slshape threadID}} to finish; it does not return any value. When it returns, it returns all
resources occupied by the thread it waited for, such as
thread\texttt{\symbol{45}}local memory and operating system structures, to the
system. }

 

\subsection{\textcolor{Chapter }{CurrentThread}}
\logpage{[ 6, 1, 3 ]}\nobreak
\hyperdef{L}{X7982BA8D782C171B}{}
{\noindent\textcolor{FuncColor}{$\triangleright$\enspace\texttt{CurrentThread({\mdseries\slshape })\index{CurrentThread@\texttt{CurrentThread}}
\label{CurrentThread}
}\hfill{\scriptsize (function)}}\\


 The \texttt{CurrentThread} function returns the thread object for the current thread. }

 

\subsection{\textcolor{Chapter }{ThreadID}}
\logpage{[ 6, 1, 4 ]}\nobreak
\hyperdef{L}{X80FA498D866028CB}{}
{\noindent\textcolor{FuncColor}{$\triangleright$\enspace\texttt{ThreadID({\mdseries\slshape thread})\index{ThreadID@\texttt{ThreadID}}
\label{ThreadID}
}\hfill{\scriptsize (function)}}\\


 The \texttt{ThreadID} function returns a numeric thread id for the given thread. The thread id of
the main thread is always 0. 
\begin{Verbatim}[commandchars=!@|,fontsize=\small,frame=single,label=Example]
  !gapprompt@gap>| !gapinput@CurrentThread();|
  <thread #0: running>
  !gapprompt@gap>| !gapinput@ThreadID(CurrentThread());|
  0
\end{Verbatim}
 }

 

\subsection{\textcolor{Chapter }{KillThread}}
\logpage{[ 6, 1, 5 ]}\nobreak
\hyperdef{L}{X787679BC7D36FAA3}{}
{\noindent\textcolor{FuncColor}{$\triangleright$\enspace\texttt{KillThread({\mdseries\slshape thread})\index{KillThread@\texttt{KillThread}}
\label{KillThread}
}\hfill{\scriptsize (function)}}\\


 The \texttt{KillThread} function terminates the given thread. Any region locks that the thread
currently holds will be unlocked. The thread can be specified as a thread
object or via its numeric id. 

 The implementation for \texttt{KillThread} is dependent on the interpreter actually executing statements. Threads
performing system calls, for example, will not be terminated until the system
call returns. Similarly, long\texttt{\symbol{45}}running kernel functions will
delay termination until the kernel function returns. 

 Use of \texttt{CALL{\textunderscore}WITH{\textunderscore}CATCH} will not prevent a thread from being terminated. If you wish to make sure that
catch handlers will be visited, use \texttt{InterruptThread} (\ref{InterruptThread}) instead. \texttt{KillThread} should be used for threads that cannot be controlled anymore in any other way
but still eat system resources. }

 

\subsection{\textcolor{Chapter }{PauseThread}}
\logpage{[ 6, 1, 6 ]}\nobreak
\hyperdef{L}{X7BC284107E56A01B}{}
{\noindent\textcolor{FuncColor}{$\triangleright$\enspace\texttt{PauseThread({\mdseries\slshape thread})\index{PauseThread@\texttt{PauseThread}}
\label{PauseThread}
}\hfill{\scriptsize (function)}}\\


 The \texttt{PauseThread} function suspends execution for the given thread. The thread can be specified
as a thread object or via its numeric id. 

 The implementation for \texttt{PauseThread} is dependent on the interpreter actually executing statements. Threads
performing system calls, for example, will not pause until the system call
returns. Similarly, long\texttt{\symbol{45}}running kernel functions will not
pause until the kernel function returns. 

 While a thread is paused, the thread that initiated the pause can access the
paused thread's thread\texttt{\symbol{45}}local region. 
\begin{Verbatim}[commandchars=!@|,fontsize=\small,frame=single,label=Example]
  !gapprompt@gap>| !gapinput@loop := function() while true do Sleep(1); od; end;;|
  !gapprompt@gap>| !gapinput@x := fail;;|
  !gapprompt@gap>| !gapinput@th := CreateThread(function() x := [1, 2, 3]; loop(); end);;|
  !gapprompt@gap>| !gapinput@PauseThread(th);|
  !gapprompt@gap>| !gapinput@x;|
  [ 1, 2, 3 ]
\end{Verbatim}
 }

 

\subsection{\textcolor{Chapter }{ResumeThread}}
\logpage{[ 6, 1, 7 ]}\nobreak
\hyperdef{L}{X80D590EC78C93F94}{}
{\noindent\textcolor{FuncColor}{$\triangleright$\enspace\texttt{ResumeThread({\mdseries\slshape thread})\index{ResumeThread@\texttt{ResumeThread}}
\label{ResumeThread}
}\hfill{\scriptsize (function)}}\\


 The \texttt{ResumeThread} function resumes execution for the given thread that was paused with \texttt{PauseThread} (\ref{PauseThread}). The thread can be specified as a thread object or via its numeric id. 

 If the thread isn't paused, \texttt{ResumeThread} is a no\texttt{\symbol{45}}op. }

 

\subsection{\textcolor{Chapter }{InterruptThread}}
\logpage{[ 6, 1, 8 ]}\nobreak
\hyperdef{L}{X8521A5777CCD0B37}{}
{\noindent\textcolor{FuncColor}{$\triangleright$\enspace\texttt{InterruptThread({\mdseries\slshape thread, interrupt})\index{InterruptThread@\texttt{InterruptThread}}
\label{InterruptThread}
}\hfill{\scriptsize (function)}}\\


 The \texttt{InterruptThread} function calls an interrupt handler for the given thread. The thread can be
specified as a thread object or via its numeric id. The interrupt is specified
as an integer between 0 and \texttt{MAX{\textunderscore}INTERRUPT} (\ref{MAXuScoreINTERRUPT}). 

 An interrupt number of zero (or an interrupt number for which no interrupt
handler has been set up with \texttt{SetInterruptHandler} (\ref{SetInterruptHandler}) will cause the thread to enter a break loop. Otherwise, the respective
interrupt handler that has been created with \texttt{SetInterruptHandler} (\ref{SetInterruptHandler}) will be called. 

 The implementation for \texttt{InterruptThread} is dependent on the interpreter actually executing statements. Threads
performing system calls, for example, will not call interrupt handlers until
the system call returns. Similarly, long\texttt{\symbol{45}}running kernel
functions will delay invocation of the interrupt handler until the kernel
function returns. }

 

\subsection{\textcolor{Chapter }{SetInterruptHandler}}
\logpage{[ 6, 1, 9 ]}\nobreak
\hyperdef{L}{X80C207B387CA888E}{}
{\noindent\textcolor{FuncColor}{$\triangleright$\enspace\texttt{SetInterruptHandler({\mdseries\slshape interrupt, handler})\index{SetInterruptHandler@\texttt{SetInterruptHandler}}
\label{SetInterruptHandler}
}\hfill{\scriptsize (function)}}\\


 The \texttt{SetInterruptHandler} function allows the programmer to set up interrupt handlers for the current
thread. The interrupt number must be in the range from 1 to \texttt{MAX{\textunderscore}INTERRUPT} (\ref{MAXuScoreINTERRUPT}) (inclusive); the handler must be a parameterless function (or \texttt{fail} to remove a handler). }

 

\subsection{\textcolor{Chapter }{NewInterruptID}}
\logpage{[ 6, 1, 10 ]}\nobreak
\hyperdef{L}{X78E5F5987BEFACE6}{}
{\noindent\textcolor{FuncColor}{$\triangleright$\enspace\texttt{NewInterruptID({\mdseries\slshape })\index{NewInterruptID@\texttt{NewInterruptID}}
\label{NewInterruptID}
}\hfill{\scriptsize (function)}}\\


 The \texttt{NewInterruptID} function returns a previously unused number (starting at 1). These numbers can
be used to globally coordinate interrupt numbers. 
\begin{Verbatim}[commandchars=!@|,fontsize=\small,frame=single,label=Example]
  !gapprompt@gap>| !gapinput@StopTaskInterrupt := NewInterruptID();|
  1
  !gapprompt@gap>| !gapinput@SetInterruptHandler(StopTaskInterrupt, StopTaskHandler);|
\end{Verbatim}
 }

 

\subsection{\textcolor{Chapter }{MAX{\textunderscore}INTERRUPT}}
\logpage{[ 6, 1, 11 ]}\nobreak
\hyperdef{L}{X863799F07D55BAE0}{}
{\noindent\textcolor{FuncColor}{$\triangleright$\enspace\texttt{MAX{\textunderscore}INTERRUPT\index{MAX{\textunderscore}INTERRUPT@\texttt{MAX{\textunderscore}INTERRUPT}}
\label{MAXuScoreINTERRUPT}
}\hfill{\scriptsize (global variable)}}\\


 The global variable \texttt{MAX{\textunderscore}INTERRUPT} is an integer containing the maximum value for the interrupt arguments to \texttt{InterruptThread} (\ref{InterruptThread}) and \texttt{SetInterruptHandler} (\ref{SetInterruptHandler}). }

 }

 }

 
\chapter{\textcolor{Chapter }{Channels}}\label{Channels}
\logpage{[ 7, 0, 0 ]}
\hyperdef{L}{X82C4DEA2841FAC0B}{}
{
  
\section{\textcolor{Chapter }{Channels}}\label{section:Channels}
\logpage{[ 7, 1, 0 ]}
\hyperdef{L}{X82C4DEA2841FAC0B}{}
{
  Channels are FIFO queues that threads can use to coordinate their activities. 

\subsection{\textcolor{Chapter }{CreateChannel}}
\logpage{[ 7, 1, 1 ]}\nobreak
\hyperdef{L}{X7E666A8B7C37ADA4}{}
{\noindent\textcolor{FuncColor}{$\triangleright$\enspace\texttt{CreateChannel({\mdseries\slshape [capacity]})\index{CreateChannel@\texttt{CreateChannel}}
\label{CreateChannel}
}\hfill{\scriptsize (function)}}\\


 \texttt{CreateChannel} returns a FIFO communication channel that can be used to exchange information
between threads. Its optional argument is a capacity (positive integer). If
insufficient resources are available to create a channel, it returns
\texttt{\symbol{45}}1. If the capacity is not a positive integer, an error
will be raised. 

 If a capacity is not provided, by default the channel can hold an indefinite
number of objects. Otherwise, attempts to store objects in the channel beyond
its capacity will block. 
\begin{Verbatim}[commandchars=!@|,fontsize=\small,frame=single,label=Example]
  !gapprompt@gap>| !gapinput@ch1:=CreateChannel();|
  <channel 0x460339c: 0 elements, 0 waiting>
  !gapprompt@gap>| !gapinput@ch2:=CreateChannel(5);|
  <channel 0x460324c: 0/5 elements, 0 waiting>
\end{Verbatim}
 }

 

\subsection{\textcolor{Chapter }{SendChannel}}
\logpage{[ 7, 1, 2 ]}\nobreak
\hyperdef{L}{X81233BCB7B1173FA}{}
{\noindent\textcolor{FuncColor}{$\triangleright$\enspace\texttt{SendChannel({\mdseries\slshape channel, obj})\index{SendChannel@\texttt{SendChannel}}
\label{SendChannel}
}\hfill{\scriptsize (function)}}\\


 \texttt{SendChannel} accepts two arguments, a channel object returned by \texttt{CreateChannel} (\ref{CreateChannel}), and an arbitrary GAP object. It stores \mbox{\texttt{\mdseries\slshape obj}} in \mbox{\texttt{\mdseries\slshape channel}}. If \mbox{\texttt{\mdseries\slshape channel}} has a finite capacity and is currently full, then \texttt{SendChannel} will block until at least one element has been removed from the channel, e.g.
using \texttt{ReceiveChannel} (\ref{ReceiveChannel}). 

 \texttt{SendChannel} performs automatic region migration for thread\texttt{\symbol{45}}local
objects. If \mbox{\texttt{\mdseries\slshape obj}} is thread\texttt{\symbol{45}}local for the current thread, it will be migrated
(along with all subobjects contained in the same region) to the receiving
thread's thread\texttt{\symbol{45}}local data space. In between sending and
receiving, \mbox{\texttt{\mdseries\slshape obj}} cannot be accessed by either thread. 

 This example demonstrates sending messages across a channel. 
\begin{Verbatim}[commandchars=!@|,fontsize=\small,frame=single,label=Example]
  !gapprompt@gap>| !gapinput@ch1 := CreateChannel();;|
  !gapprompt@gap>| !gapinput@SendChannel(ch1,1);|
  !gapprompt@gap>| !gapinput@ch1;|
  <channel 0x460339c: 1 elements, 0 waiting>
  !gapprompt@gap>| !gapinput@ReceiveChannel(ch1);|
  1
  !gapprompt@gap>| !gapinput@ch1;|
  <channel 0x460339c: 0 elements, 0 waiting>
\end{Verbatim}
 \texttt{Sleep} in the following example is used to demonstrate blocking. 
\begin{Verbatim}[commandchars=!@|,fontsize=\small,frame=single,label=Example]
  !gapprompt@gap>| !gapinput@ch2 := CreateChannel(5);;|
  !gapprompt@gap>| !gapinput@ch3 := CreateChannel();;|
  !gapprompt@gap>| !gapinput@for i in [1..5] do SendChannel(ch2,i); od;|
  !gapprompt@gap>| !gapinput@ch2;|
  <channel 0x460324c: 5/5 elements, 0 waiting>
  !gapprompt@gap>| !gapinput@t:=CreateThread(|
  !gapprompt@>| !gapinput@function()|
  !gapprompt@>| !gapinput@local x;|
  !gapprompt@>| !gapinput@Sleep(10);|
  !gapprompt@>| !gapinput@x:=ReceiveChannel(ch2);|
  !gapprompt@>| !gapinput@Sleep(10);|
  !gapprompt@>| !gapinput@SendChannel(ch3,x);|
  !gapprompt@>| !gapinput@Print("Thread finished\n");|
  !gapprompt@>| !gapinput@end);;|
  !gapprompt@>| !gapinput@SendChannel(ch2,3); # this blocks until the thread reads from ch2|
  !gapprompt@gap>| !gapinput@ReceiveChannel(ch3); # the thread is blocked until we read from ch3|
  1
  Thread finished
  !gapprompt@gap>| !gapinput@WaitThread(t);|
\end{Verbatim}
 }

 

\subsection{\textcolor{Chapter }{TransmitChannel}}
\logpage{[ 7, 1, 3 ]}\nobreak
\hyperdef{L}{X8004E54A7D91929C}{}
{\noindent\textcolor{FuncColor}{$\triangleright$\enspace\texttt{TransmitChannel({\mdseries\slshape channel, obj})\index{TransmitChannel@\texttt{TransmitChannel}}
\label{TransmitChannel}
}\hfill{\scriptsize (function)}}\\


 \texttt{TransmitChannel} is identical to \texttt{SendChannel} (\ref{SendChannel}), except that it does not perform automatic region migration of
thread\texttt{\symbol{45}}local objects. 
\begin{Verbatim}[commandchars=!@|,fontsize=\small,frame=single,label=Example]
  !gapprompt@gap>| !gapinput@ch := CreateChannel(5);;|
  !gapprompt@gap>| !gapinput@l := [ 1, 2, 3];;|
  !gapprompt@gap>| !gapinput@original_region := RegionOf(l);;|
  !gapprompt@gap>| !gapinput@SendChannel(ch, l);|
  !gapprompt@gap>| !gapinput@WaitThread(CreateThread(function()|
  !gapprompt@>| !gapinput@     local ob; ob := ReceiveChannel(ch);|
  !gapprompt@>| !gapinput@     Display(RegionOf(ob) = original_region);|
  !gapprompt@>| !gapinput@   end));|
  false
  !gapprompt@gap>| !gapinput@l := [ 1, 2, 3];;|
  !gapprompt@gap>| !gapinput@original_region := RegionOf(l);;|
  !gapprompt@gap>| !gapinput@TransmitChannel(ch, l);|
  !gapprompt@gap>| !gapinput@WaitThread(CreateThread(function()|
  !gapprompt@>| !gapinput@     local ob; ob := ReceiveChannel(ch);|
  !gapprompt@>| !gapinput@     Display(RegionOf(ob) = original_region);|
  !gapprompt@>| !gapinput@   end));|
  true
\end{Verbatim}
 }

 

\subsection{\textcolor{Chapter }{TrySendChannel}}
\logpage{[ 7, 1, 4 ]}\nobreak
\hyperdef{L}{X7AEFF25D8143DA77}{}
{\noindent\textcolor{FuncColor}{$\triangleright$\enspace\texttt{TrySendChannel({\mdseries\slshape channel, obj})\index{TrySendChannel@\texttt{TrySendChannel}}
\label{TrySendChannel}
}\hfill{\scriptsize (function)}}\\


 \texttt{TrySendChannel} is identical to \texttt{SendChannel} (\ref{SendChannel}), except that it returns if the channel is full instead of blocking. It
returns true if the send was successful and false otherwise. 
\begin{Verbatim}[commandchars=!@|,fontsize=\small,frame=single,label=Example]
  !gapprompt@gap>| !gapinput@ch := CreateChannel(1);;|
  !gapprompt@gap>| !gapinput@TrySendChannel(ch, 99);|
  true
  !gapprompt@gap>| !gapinput@TrySendChannel(ch, 99);|
  false
\end{Verbatim}
 }

 

\subsection{\textcolor{Chapter }{TryTransmitChannel}}
\logpage{[ 7, 1, 5 ]}\nobreak
\hyperdef{L}{X7BB33C037A81BCF6}{}
{\noindent\textcolor{FuncColor}{$\triangleright$\enspace\texttt{TryTransmitChannel({\mdseries\slshape channel, obj})\index{TryTransmitChannel@\texttt{TryTransmitChannel}}
\label{TryTransmitChannel}
}\hfill{\scriptsize (function)}}\\


 \texttt{TryTransmitChannel} is identical to \texttt{TrySendChannel} (\ref{TrySendChannel}), except that it does not perform automatic region migration of
thread\texttt{\symbol{45}}local objects. }

 

\subsection{\textcolor{Chapter }{ReceiveChannel}}
\logpage{[ 7, 1, 6 ]}\nobreak
\hyperdef{L}{X7FEABC588101CCE7}{}
{\noindent\textcolor{FuncColor}{$\triangleright$\enspace\texttt{ReceiveChannel({\mdseries\slshape channel})\index{ReceiveChannel@\texttt{ReceiveChannel}}
\label{ReceiveChannel}
}\hfill{\scriptsize (function)}}\\


 \texttt{ReceiveChannel} is used to retrieve elements from a channel. If \mbox{\texttt{\mdseries\slshape channel}} is empty, the call will block until an element has been added to the channel
via \texttt{SendChannel} (\ref{SendChannel}) or a similar primitive. 

 See \texttt{SendChannel} (\ref{SendChannel}) for an example. }

 

\subsection{\textcolor{Chapter }{TryReceiveChannel}}
\logpage{[ 7, 1, 7 ]}\nobreak
\hyperdef{L}{X788785E380C10A01}{}
{\noindent\textcolor{FuncColor}{$\triangleright$\enspace\texttt{TryReceiveChannel({\mdseries\slshape channel, default})\index{TryReceiveChannel@\texttt{TryReceiveChannel}}
\label{TryReceiveChannel}
}\hfill{\scriptsize (function)}}\\


 \texttt{TryReceiveChannel}, like \texttt{ReceiveChannel} (\ref{ReceiveChannel}), attempts to retrieve an object from \mbox{\texttt{\mdseries\slshape channel}}. If it does not succeed, however, it will return \mbox{\texttt{\mdseries\slshape default}} rather than blocking. 
\begin{Verbatim}[commandchars=!@|,fontsize=\small,frame=single,label=Example]
  !gapprompt@gap>| !gapinput@ch := CreateChannel();;|
  !gapprompt@gap>| !gapinput@SendChannel(ch, 99);|
  !gapprompt@gap>| !gapinput@TryReceiveChannel(ch, fail);|
  99
  !gapprompt@gap>| !gapinput@TryReceiveChannel(ch, fail);|
  fail
\end{Verbatim}
 }

 

\subsection{\textcolor{Chapter }{MultiSendChannel}}
\logpage{[ 7, 1, 8 ]}\nobreak
\hyperdef{L}{X7A5E26F5799A17B4}{}
{\noindent\textcolor{FuncColor}{$\triangleright$\enspace\texttt{MultiSendChannel({\mdseries\slshape channel, list})\index{MultiSendChannel@\texttt{MultiSendChannel}}
\label{MultiSendChannel}
}\hfill{\scriptsize (function)}}\\


 \texttt{MultiSendChannel} allows the sending of all the objects contained in the list \mbox{\texttt{\mdseries\slshape list}} to \mbox{\texttt{\mdseries\slshape channel}} as a single operation. The list must be dense and is not modified by the call.
The function will send elements starting at index 1 until all elements have
been sent. If a channel with finite capacity is full, then the operation will
block until all elements can be sent. 

 The operation is designed to be more efficient than sending all elements
individually via \texttt{SendChannel} (\ref{SendChannel}) by minimizing potentially expensive concurrency operations. 

 See \texttt{MultiReceiveChannel} (\ref{MultiReceiveChannel}) for an example. }

 

\subsection{\textcolor{Chapter }{TryMultiSendChannel}}
\logpage{[ 7, 1, 9 ]}\nobreak
\hyperdef{L}{X85B8D09B802D2122}{}
{\noindent\textcolor{FuncColor}{$\triangleright$\enspace\texttt{TryMultiSendChannel({\mdseries\slshape channel, list})\index{TryMultiSendChannel@\texttt{TryMultiSendChannel}}
\label{TryMultiSendChannel}
}\hfill{\scriptsize (function)}}\\


 \texttt{TryMultiSendChannel} operates like \texttt{MultiSendChannel} (\ref{MultiSendChannel}), except that it returns rather than blocking if it cannot send any more
elements if the channel is full. It returns the number of elements it has
sent. If \mbox{\texttt{\mdseries\slshape channel}} does not have finite capacity, \texttt{TryMultiSendChannel} will always send all elements in the list. }

 

\subsection{\textcolor{Chapter }{MultiReceiveChannel}}
\logpage{[ 7, 1, 10 ]}\nobreak
\hyperdef{L}{X807D4263798D775E}{}
{\noindent\textcolor{FuncColor}{$\triangleright$\enspace\texttt{MultiReceiveChannel({\mdseries\slshape channel, amount})\index{MultiReceiveChannel@\texttt{MultiReceiveChannel}}
\label{MultiReceiveChannel}
}\hfill{\scriptsize (function)}}\\


 \texttt{MultiReceiveChannel} is the receiving counterpart to \texttt{MultiSendChannel} (\ref{MultiSendChannel}).It will try to receive up to \mbox{\texttt{\mdseries\slshape amount}} objects from \mbox{\texttt{\mdseries\slshape channel}}. If the channel contains less than \mbox{\texttt{\mdseries\slshape amount}} objects, it will return rather than blocking. 

 The function returns a list of all the objects received. 
\begin{Verbatim}[commandchars=!@|,fontsize=\small,frame=single,label=Example]
  !gapprompt@gap>| !gapinput@ch:=CreateChannel();;|
  !gapprompt@gap>| !gapinput@MultiSendChannel(ch, [1, 2, 3, 4, 5, 6, 7, 8, 9, 10]);|
  !gapprompt@gap>| !gapinput@MultiReceiveChannel(ch,7);|
  [ 1, 2, 3, 4, 5, 6, 7 ]
  !gapprompt@gap>| !gapinput@MultiReceiveChannel(ch,7);|
  [ 8, 9, 10 ]
  !gapprompt@gap>| !gapinput@MultiReceiveChannel(ch,7);|
  [  ]
\end{Verbatim}
 }

 

\subsection{\textcolor{Chapter }{ReceiveAnyChannel}}
\logpage{[ 7, 1, 11 ]}\nobreak
\hyperdef{L}{X8355AFCC8530B2CA}{}
{\noindent\textcolor{FuncColor}{$\triangleright$\enspace\texttt{ReceiveAnyChannel({\mdseries\slshape channel{\textunderscore}1, ..., channel{\textunderscore}n})\index{ReceiveAnyChannel@\texttt{ReceiveAnyChannel}}
\label{ReceiveAnyChannel}
}\hfill{\scriptsize (function)}}\\
\noindent\textcolor{FuncColor}{$\triangleright$\enspace\texttt{ReceiveAnyChannel({\mdseries\slshape channel{\textunderscore}list})\index{ReceiveAnyChannel@\texttt{ReceiveAnyChannel}!for a list of channels}
\label{ReceiveAnyChannel:for a list of channels}
}\hfill{\scriptsize (function)}}\\


 \texttt{ReceiveAnyChannel} is a multiplexing variant of\texttt{ReceiveChannel} (\ref{ReceiveChannel}). It blocks until at least one of the channels provided contains an object. It
will then retrieve that object from the channel and return it. 
\begin{Verbatim}[commandchars=!@|,fontsize=\small,frame=single,label=Example]
  !gapprompt@gap>| !gapinput@ch1 := CreateChannel();;|
  !gapprompt@gap>| !gapinput@ch2 := CreateChannel();;|
  !gapprompt@gap>| !gapinput@SendChannel(ch2, [1, 2, 3]);;|
  !gapprompt@gap>| !gapinput@ReceiveAnyChannel(ch1, ch2);|
  [ 1, 2, 3 ]
\end{Verbatim}
 }

 

\subsection{\textcolor{Chapter }{ReceiveAnyChannelWithIndex}}
\logpage{[ 7, 1, 12 ]}\nobreak
\hyperdef{L}{X846C7B3F7BF7CA0B}{}
{\noindent\textcolor{FuncColor}{$\triangleright$\enspace\texttt{ReceiveAnyChannelWithIndex({\mdseries\slshape channel{\textunderscore}1, ..., channel{\textunderscore}n})\index{ReceiveAnyChannelWithIndex@\texttt{ReceiveAnyChannelWithIndex}}
\label{ReceiveAnyChannelWithIndex}
}\hfill{\scriptsize (function)}}\\
\noindent\textcolor{FuncColor}{$\triangleright$\enspace\texttt{ReceiveAnyChannelWithIndex({\mdseries\slshape channel{\textunderscore}list})\index{ReceiveAnyChannelWithIndex@\texttt{ReceiveAnyChannelWithIndex}!for a list of channels}
\label{ReceiveAnyChannelWithIndex:for a list of channels}
}\hfill{\scriptsize (function)}}\\


 \texttt{ReceiveAnyChannelWithIndex} works like \texttt{ReceiveAnyChannel} (\ref{ReceiveAnyChannel}), except that it returns a list with two elements, the first being the object
being received, the second being the number of the channel from which the
object has been retrieved. 
\begin{Verbatim}[commandchars=!@|,fontsize=\small,frame=single,label=Example]
  !gapprompt@gap>| !gapinput@ch1 := CreateChannel();;|
  !gapprompt@gap>| !gapinput@ch2 := CreateChannel();;|
  !gapprompt@gap>| !gapinput@SendChannel(ch2, [1, 2, 3]);;|
  !gapprompt@gap>| !gapinput@ReceiveAnyChannelWithIndex(ch1, ch2);|
  [ [ 1, 2, 3 ], 2 ]
\end{Verbatim}
 }

 

\subsection{\textcolor{Chapter }{TallyChannel}}
\logpage{[ 7, 1, 13 ]}\nobreak
\hyperdef{L}{X81A731667C33E150}{}
{\noindent\textcolor{FuncColor}{$\triangleright$\enspace\texttt{TallyChannel({\mdseries\slshape channel})\index{TallyChannel@\texttt{TallyChannel}}
\label{TallyChannel}
}\hfill{\scriptsize (function)}}\\


 \texttt{TallyChannel} returns the number of objects that a channel contains. This number can
increase or decrease, as data is sent to or received from this channel. Send
operations will only ever increase and receive operations will only ever
decrease this count. Thus, if there is only one thread receiving data from the
channel, it can use the result as a lower bound for the number of elements
that will be available in the channel. 
\begin{Verbatim}[commandchars=!@|,fontsize=\small,frame=single,label=Example]
  !gapprompt@gap>| !gapinput@ch := CreateChannel();;|
  !gapprompt@gap>| !gapinput@SendChannel(ch, 2);|
  !gapprompt@gap>| !gapinput@SendChannel(ch, 3);|
  !gapprompt@gap>| !gapinput@SendChannel(ch, 5);|
  !gapprompt@gap>| !gapinput@TallyChannel(ch);|
  3
\end{Verbatim}
 }

 

\subsection{\textcolor{Chapter }{InspectChannel}}
\logpage{[ 7, 1, 14 ]}\nobreak
\hyperdef{L}{X7AF2B2ED8552E33D}{}
{\noindent\textcolor{FuncColor}{$\triangleright$\enspace\texttt{InspectChannel({\mdseries\slshape channel})\index{InspectChannel@\texttt{InspectChannel}}
\label{InspectChannel}
}\hfill{\scriptsize (function)}}\\


 \texttt{InspectChannel} returns a list of the objects that a channel contains. Note that objects that
are not in the shared, public, or read\texttt{\symbol{45}}only region will be
temporarily stored in the so\texttt{\symbol{45}}called limbo region while in
transit and will be inaccessible through normal means until they have been
received. 
\begin{Verbatim}[commandchars=!@|,fontsize=\small,frame=single,label=Example]
  !gapprompt@gap>| !gapinput@ch := CreateChannel();;|
  !gapprompt@gap>| !gapinput@SendChannel(ch, 2);|
  !gapprompt@gap>| !gapinput@SendChannel(ch, 3);|
  !gapprompt@gap>| !gapinput@SendChannel(ch, 5);|
  !gapprompt@gap>| !gapinput@InspectChannel(ch);|
  [ 2, 3, 5 ]
\end{Verbatim}
 This function is primarily intended for debugging purposes. }

 }

 }

 
\chapter{\textcolor{Chapter }{Semaphores}}\label{Semaphores}
\logpage{[ 8, 0, 0 ]}
\hyperdef{L}{X83F63F0F7827767E}{}
{
  
\section{\textcolor{Chapter }{Semaphores}}\label{section:Semaphores}
\logpage{[ 8, 1, 0 ]}
\hyperdef{L}{X83F63F0F7827767E}{}
{
  Semaphores are synchronized counters; they can also be used to simulate locks. 

\subsection{\textcolor{Chapter }{CreateSemaphore}}
\logpage{[ 8, 1, 1 ]}\nobreak
\hyperdef{L}{X7AC5500483465AE3}{}
{\noindent\textcolor{FuncColor}{$\triangleright$\enspace\texttt{CreateSemaphore({\mdseries\slshape [value]})\index{CreateSemaphore@\texttt{CreateSemaphore}}
\label{CreateSemaphore}
}\hfill{\scriptsize (function)}}\\


 The function \texttt{CreateSemaphore} takes an optional argument, which defaults to zero. It is the counter with
which the semaphore is initialized. 
\begin{Verbatim}[commandchars=!@|,fontsize=\small,frame=single,label=Example]
  !gapprompt@gap>| !gapinput@sem := CreateSemaphore(1);|
  <semaphore 0x1108e81c0: count = 1>
\end{Verbatim}
 }

 

\subsection{\textcolor{Chapter }{WaitSemaphore}}
\logpage{[ 8, 1, 2 ]}\nobreak
\hyperdef{L}{X834E47327A3FC5A2}{}
{\noindent\textcolor{FuncColor}{$\triangleright$\enspace\texttt{WaitSemaphore({\mdseries\slshape sem})\index{WaitSemaphore@\texttt{WaitSemaphore}}
\label{WaitSemaphore}
}\hfill{\scriptsize (function)}}\\


 \texttt{WaitSemaphore} receives a previously created semaphore as its argument. If the semaphore's
counter is greater than zero, it decrements the counter and returns; if the
counter is zero, it waits until another thread increases it via \texttt{SignalSemaphore} (\ref{SignalSemaphore}), then decrements the counter and returns. 
\begin{Verbatim}[commandchars=!@|,fontsize=\small,frame=single,label=Example]
  !gapprompt@gap>| !gapinput@sem := CreateSemaphore(1);|
  <semaphore 0x1108e81c0: count = 1>
  !gapprompt@gap>| !gapinput@WaitSemaphore(sem);|
  !gapprompt@gap>| !gapinput@sem;|
  <semaphore 0x1108e81c0: count = 0>
\end{Verbatim}
 }

 

\subsection{\textcolor{Chapter }{SignalSemaphore}}
\logpage{[ 8, 1, 3 ]}\nobreak
\hyperdef{L}{X7F0C1F1F8540CF8C}{}
{\noindent\textcolor{FuncColor}{$\triangleright$\enspace\texttt{SignalSemaphore({\mdseries\slshape sem})\index{SignalSemaphore@\texttt{SignalSemaphore}}
\label{SignalSemaphore}
}\hfill{\scriptsize (function)}}\\


 \texttt{SignalSemaphore} receives a previously created semaphore as its argument. It increments the
semaphore's counter and returns. 
\begin{Verbatim}[commandchars=!@|,fontsize=\small,frame=single,label=Example]
  !gapprompt@gap>| !gapinput@sem := CreateSemaphore(1);|
  <semaphore 0x1108e81c0: count = 1>
  !gapprompt@gap>| !gapinput@WaitSemaphore(sem);|
  !gapprompt@gap>| !gapinput@sem;|
  <semaphore 0x1108e81c0: count = 0>
  !gapprompt@gap>| !gapinput@SignalSemaphore(sem);|
  !gapprompt@gap>| !gapinput@sem;|
  <semaphore 0x1108e81c0: count = 1>
\end{Verbatim}
 }

 
\subsection{\textcolor{Chapter }{Simulating locks}}\label{Simulating locks}
\logpage{[ 8, 1, 4 ]}
\hyperdef{L}{X83D868417F0069F0}{}
{
  In order to use semaphores to simulate locks, create a semaphore with an
initial value of 1. \texttt{WaitSemaphore} (\ref{WaitSemaphore}) is then equivalent to a lock operation, \texttt{SignalSemaphore} (\ref{SignalSemaphore}) is equivalent to an unlock operation. }

 }

 }

 
\chapter{\textcolor{Chapter }{Synchronization variables}}\label{Synchronization variables}
\logpage{[ 9, 0, 0 ]}
\hyperdef{L}{X84583568784B622A}{}
{
  
\section{\textcolor{Chapter }{Synchronization variables}}\label{section:Synchronization variables}
\logpage{[ 9, 1, 0 ]}
\hyperdef{L}{X84583568784B622A}{}
{
  Synchronization variables (also often called dataflow variables in the
literature) are variables that can be written only once; attempts to read the
variable block until it has been written to. 

 Synchronization variables are created with \texttt{CreateSyncVar} (\ref{CreateSyncVar}), written with \texttt{SyncWrite} (\ref{SyncWrite}) and read with \texttt{SyncRead} (\ref{SyncRead}). 
\begin{Verbatim}[commandchars=!@|,fontsize=\small,frame=single,label=Example]
  !gapprompt@gap>| !gapinput@sv := CreateSyncVar();;|
  !gapprompt@gap>| !gapinput@RunAsyncTask(function()|
  !gapprompt@>| !gapinput@     Sleep(10);|
  !gapprompt@>| !gapinput@     SyncWrite(sv, MakeImmutable([1, 2, 3]));|
  !gapprompt@>| !gapinput@   end);;|
  !gapprompt@gap>| !gapinput@SyncRead(sv);|
  [ 1, 2, 3 ]
\end{Verbatim}
 

\subsection{\textcolor{Chapter }{CreateSyncVar}}
\logpage{[ 9, 1, 1 ]}\nobreak
\hyperdef{L}{X7B55EA0E8672087C}{}
{\noindent\textcolor{FuncColor}{$\triangleright$\enspace\texttt{CreateSyncVar({\mdseries\slshape })\index{CreateSyncVar@\texttt{CreateSyncVar}}
\label{CreateSyncVar}
}\hfill{\scriptsize (function)}}\\


 The function \texttt{CreateSyncVar} takes no arguments. It returns a new synchronization variable. There is no
need to deallocate it; the garbage collector will free the memory and all
related resources when it is no longer accessible. }

 

\subsection{\textcolor{Chapter }{SyncWrite}}
\logpage{[ 9, 1, 2 ]}\nobreak
\hyperdef{L}{X85E910BF7BBF3270}{}
{\noindent\textcolor{FuncColor}{$\triangleright$\enspace\texttt{SyncWrite({\mdseries\slshape syncvar, obj})\index{SyncWrite@\texttt{SyncWrite}}
\label{SyncWrite}
}\hfill{\scriptsize (function)}}\\


 \texttt{SyncWrite} attempts to assign the value \mbox{\texttt{\mdseries\slshape obj}} to \mbox{\texttt{\mdseries\slshape syncvar}}. If \mbox{\texttt{\mdseries\slshape syncvar}} has been previously assigned a value, the call will fail with a runtime error;
otherwise, \mbox{\texttt{\mdseries\slshape obj}} will be assigned to \mbox{\texttt{\mdseries\slshape syncvar}}. 

 In order to make sure that the recipient can read the result, the \mbox{\texttt{\mdseries\slshape obj}} argument should not be a thread\texttt{\symbol{45}}local object; it should be
public, read\texttt{\symbol{45}}only, or shared. }

 

\subsection{\textcolor{Chapter }{SyncRead}}
\logpage{[ 9, 1, 3 ]}\nobreak
\hyperdef{L}{X7B42B29B8441B4F5}{}
{\noindent\textcolor{FuncColor}{$\triangleright$\enspace\texttt{SyncRead({\mdseries\slshape syncvar})\index{SyncRead@\texttt{SyncRead}}
\label{SyncRead}
}\hfill{\scriptsize (function)}}\\


 \texttt{SyncRead} reads the value previously assigned to \mbox{\texttt{\mdseries\slshape syncvar}} with \texttt{SyncWrite} (\ref{SyncWrite}). If no value has been assigned yet, it blocks. It returns the assigned value. }

 }

 }

 
\chapter{\textcolor{Chapter }{Serialization support}}\label{Serialization support}
\logpage{[ 10, 0, 0 ]}
\hyperdef{L}{X806C386986781A27}{}
{
  
\section{\textcolor{Chapter }{Serialization support}}\label{section:Serialization support}
\logpage{[ 10, 1, 0 ]}
\hyperdef{L}{X806C386986781A27}{}
{
  HPC\texttt{\symbol{45}}GAP has support to serialize most GAP data. While
functions in particular cannot be serialized, it is possible to serialize all
primitive types (booleans, integers, cyclotomics, permutations, floats, etc.)
as well as all lists and records. 

 Custom serialization support can be written for data objects, positional
objects, and component objects; serialization of compressed vectors is already
supported by the standard library. 

\subsection{\textcolor{Chapter }{SerializeToNativeString}}
\logpage{[ 10, 1, 1 ]}\nobreak
\hyperdef{L}{X787953287A989F91}{}
{\noindent\textcolor{FuncColor}{$\triangleright$\enspace\texttt{SerializeToNativeString({\mdseries\slshape obj})\index{SerializeToNativeString@\texttt{SerializeToNativeString}}
\label{SerializeToNativeString}
}\hfill{\scriptsize (function)}}\\


 \texttt{SerializeToNativeString} takes the object passed as an argument and turns it into a string, from which
a copy of the original can be extracted using \texttt{DeserializeNativeString} (\ref{DeserializeNativeString}). }

 

\subsection{\textcolor{Chapter }{DeserializeNativeString}}
\logpage{[ 10, 1, 2 ]}\nobreak
\hyperdef{L}{X80FCBCCF785AA385}{}
{\noindent\textcolor{FuncColor}{$\triangleright$\enspace\texttt{DeserializeNativeString({\mdseries\slshape str})\index{DeserializeNativeString@\texttt{DeserializeNativeString}}
\label{DeserializeNativeString}
}\hfill{\scriptsize (function)}}\\


 \texttt{DeserializeNativeString} reverts the serialization process. 

 Example: 
\begin{Verbatim}[commandchars=!@|,fontsize=\small,frame=single,label=Example]
  !gapprompt@gap>| !gapinput@DeserializeNativeString(SerializeToNativeString([1,2,3]));|
  [ 1, 2, 3 ]
\end{Verbatim}
 }

 

\subsection{\textcolor{Chapter }{InstallTypeSerializationTag}}
\logpage{[ 10, 1, 3 ]}\nobreak
\hyperdef{L}{X7E8286777A438AB5}{}
{\noindent\textcolor{FuncColor}{$\triangleright$\enspace\texttt{InstallTypeSerializationTag({\mdseries\slshape type, tag})\index{InstallTypeSerializationTag@\texttt{InstallTypeSerializationTag}}
\label{InstallTypeSerializationTag}
}\hfill{\scriptsize (function)}}\\


 \texttt{InstallTypeSerializationTag} allows the serialization of data objects, positional objects, and component
objects. The value of \mbox{\texttt{\mdseries\slshape tag}} must be unique for each type; it can be a string or integer.
Non\texttt{\symbol{45}}negative integers are reserved for use by the standard
library; users should use negative integers or strings instead. 

 Objects of such a type are serialized in a straightforward way: During
serialization, data objects are converted into byte streams, positional
objects into lists, and component objects into records. These objects are then
serialized along with their tags; deserialization uses the type corresponding
to the tag in conjunction with \texttt{Objectify} (\textbf{Reference: Objectify}) to reconstruct a copy of the original object. 

 Note that this functionality may be inadequate for objects that have complex
data structures attached that are not meant to be replicated. The following
alternative is meant for such objects. }

 

\subsection{\textcolor{Chapter }{InstallSerializer}}
\logpage{[ 10, 1, 4 ]}\nobreak
\hyperdef{L}{X84E98E728481BE0F}{}
{\noindent\textcolor{FuncColor}{$\triangleright$\enspace\texttt{InstallSerializer({\mdseries\slshape description, filters, method})\index{InstallSerializer@\texttt{InstallSerializer}}
\label{InstallSerializer}
}\hfill{\scriptsize (function)}}\\


 The more general \texttt{InstallSerializer} allows for arbitrarily complex serialization code. It installs \mbox{\texttt{\mdseries\slshape method}} as the method to serialize objects matching \mbox{\texttt{\mdseries\slshape filters}}; \mbox{\texttt{\mdseries\slshape description}} has the same role as for \texttt{InstallMethod} (\textbf{Reference: InstallMethod}). 

 The method must return a plain list matching a specific format. The first
element must be a non\texttt{\symbol{45}}negative integer, the second must be
a string descriptor that is unique to the serializer; these can then be
followed by an arbitrary number of arguments. 

 As many of the arguments (starting with the third element of the list) as
specified by the first element of the list will be converted from their object
representation into a serializable representation. Data objects will be
converted into untyped data objects, positional objects will be converted into
plain lists, and component objects into records. Conversion will not modify
the objects in place, but work on copies. The remaining arguments will remain
untouched. 

 Upon deserialization, these arguments will be passed to a function specified
by the second element of the list. 

 Example: 
\begin{Verbatim}[commandchars=!@|,fontsize=\small,frame=single,label=Example]
  InstallSerializer("8-bit vectors", [ Is8BitVectorRep ], function(obj)
    return [1, "Vec8Bit", obj, Q_VEC8BIT(obj), IS_MUTABLE_OBJ(obj)];
  end);
\end{Verbatim}
 Here, \texttt{obj} will be converted into its underlying representation, while the remaining
arguments are left alone. \texttt{"Vec8Bit"} is the name that is used to look up the deserializer function. }

 

\subsection{\textcolor{Chapter }{InstallDeserializer}}
\logpage{[ 10, 1, 5 ]}\nobreak
\hyperdef{L}{X7A9117BE7BE87CD6}{}
{\noindent\textcolor{FuncColor}{$\triangleright$\enspace\texttt{InstallDeserializer({\mdseries\slshape descriptor, func})\index{InstallDeserializer@\texttt{InstallDeserializer}}
\label{InstallDeserializer}
}\hfill{\scriptsize (function)}}\\


 The \mbox{\texttt{\mdseries\slshape descriptor}} value must be the same as the second element of the list returned by the
serializer; \mbox{\texttt{\mdseries\slshape func}} must be a function that takes as many arguments as there were arguments after
the second element of that list. For deserialization, this function is invoked
and needs to return the deserialized object constructed from the arguments. 

 Example: 
\begin{Verbatim}[commandchars=!@|,fontsize=\small,frame=single,label=Example]
  InstallDeserializer("Vec8Bit", function(obj, q, mut)
    SET_TYPE_OBJ(obj, TYPE_VEC8BIT(q, mut));
    return obj;
  end);
\end{Verbatim}
 Here, the untyped \texttt{obj} that was passed to the deserializer needs to be given the correct type, which
is calculated from \texttt{q} and \texttt{mut}. }

 }

 }

 
\chapter{\textcolor{Chapter }{Low\texttt{\symbol{45}}level functionality}}\label{Low-level functionality}
\logpage{[ 11, 0, 0 ]}
\hyperdef{L}{X7B27C7E98204353D}{}
{
  The functionality described in this section should only be used by experts,
and even by those only with caution (especially the parts that relate to the
memory model). 

 Not only is it possible to crash or hang the GAP kernel, it can happen in ways
that are very difficult to reproduce, leading to software defects that are
discovered only long after deployment of a package and then become difficult
to correct. 

 The performance benefit of using these primitives is generally minimal; while
concurrency can induce some overhead, the benefit from micromanaging
concurrency in an interpreted language such as GAP is likely to be small. 

 These low\texttt{\symbol{45}}level primitives exist primarily for the benefit
of kernel programmers; it allows them to prototype new kernel functionality in
GAP before implementing it in C. 
\section{\textcolor{Chapter }{Explicit lock and unlock primitives}}\label{Explicit lock and unlock primitives}
\logpage{[ 11, 1, 0 ]}
\hyperdef{L}{X864926CD80307F32}{}
{
  The \texttt{LOCK} (\ref{LOCK}) operation combined with \texttt{UNLOCK} (\ref{UNLOCK}) is a low\texttt{\symbol{45}}level interface for the functionality of the
statement. 

\subsection{\textcolor{Chapter }{LOCK}}
\logpage{[ 11, 1, 1 ]}\nobreak
\hyperdef{L}{X828E21CD78EFE07A}{}
{\noindent\textcolor{FuncColor}{$\triangleright$\enspace\texttt{LOCK({\mdseries\slshape [arg{\textunderscore}1, ..., arg{\textunderscore}n]})\index{LOCK@\texttt{LOCK}}
\label{LOCK}
}\hfill{\scriptsize (function)}}\\


 \texttt{LOCK} takes zero or more pairs of parameters, where each is either an object or a
boolean value. If an argument is an object, the region containing it will be
locked. If an argument is the boolean value \texttt{false}, all subsequent locks will be read locks; if it is the boolean value \texttt{true}, all subsequent locks will be write locks. If the first argument is not a
boolean value, all locks until the first boolean value will be write locks. 

 Locks are managed internally as a stack of locked regions; \texttt{LOCK} returns an integer indicating a pointer to the top of the stack; this integer
is used later by the \texttt{UNLOCK} (\ref{UNLOCK}) operation to unlock locks on the stack up to that position. If \texttt{LOCK} should fail for some reason, it will return \texttt{fail}. 

 Calling \texttt{LOCK} with no parameters returns the current lock stack pointer. }

 

\subsection{\textcolor{Chapter }{TRYLOCK}}
\logpage{[ 11, 1, 2 ]}\nobreak
\hyperdef{L}{X7A7256AA7D7E3C4B}{}
{\noindent\textcolor{FuncColor}{$\triangleright$\enspace\texttt{TRYLOCK({\mdseries\slshape [arg{\textunderscore}1, ..., arg{\textunderscore}n]})\index{TRYLOCK@\texttt{TRYLOCK}}
\label{TRYLOCK}
}\hfill{\scriptsize (function)}}\\


 \texttt{TRYLOCK} works similarly to \texttt{LOCK} (\ref{LOCK}). If it cannot acquire all region locks, it returns \texttt{fail} and does not lock any regions. Otherwise, its semantics are identical to \texttt{LOCK} (\ref{LOCK}). }

 

\subsection{\textcolor{Chapter }{UNLOCK}}
\logpage{[ 11, 1, 3 ]}\nobreak
\hyperdef{L}{X7EBEFBC6875F149A}{}
{\noindent\textcolor{FuncColor}{$\triangleright$\enspace\texttt{UNLOCK({\mdseries\slshape stackpos})\index{UNLOCK@\texttt{UNLOCK}}
\label{UNLOCK}
}\hfill{\scriptsize (function)}}\\


 \texttt{UNLOCK} unlocks all regions on the stack at \mbox{\texttt{\mdseries\slshape stackpos}} or higher and sets the stack pointer to \mbox{\texttt{\mdseries\slshape stackpos}}. 
\begin{Verbatim}[commandchars=!@|,fontsize=\small,frame=single,label=Example]
  !gapprompt@gap>| !gapinput@l1 := ShareObj([1,2,3]);;|
  !gapprompt@gap>| !gapinput@l2 := ShareObj([4,5,6]);;|
  !gapprompt@gap>| !gapinput@p := LOCK(l1);|
  0
  !gapprompt@gap>| !gapinput@LOCK(l2);|
  1
  !gapprompt@gap>| !gapinput@UNLOCK(p); # unlock both RegionOf(l1) and RegionOf(l2)|
  !gapprompt@gap>| !gapinput@LOCK(); # current stack pointer|
  0
\end{Verbatim}
 }

 }

 
\section{\textcolor{Chapter }{Hash locks}}\label{Hash locks}
\logpage{[ 11, 2, 0 ]}
\hyperdef{L}{X83B01F7D7DC9FCD0}{}
{
  HPC\texttt{\symbol{45}}GAP supports \emph{hash locks}; internally, the kernel maintains a fixed size array of locks; objects are
mapped to a lock via hash function. The hash function is based on the object
reference, not its contents (except for short integers and finite field
elements). 
\begin{Verbatim}[commandchars=!@|,fontsize=\small,frame=single,label=Example]
  !gapprompt@gap>| !gapinput@l := [ 1, 2, 3];;|
  !gapprompt@gap>| !gapinput@f := l -> Sum(l);;|
  !gapprompt@gap>| !gapinput@HASH_LOCK(l);   # lock 'l'|
  !gapprompt@gap>| !gapinput@f(l);           # do something with 'l'|
  6
  !gapprompt@gap>| !gapinput@HASH_UNLOCK(l); # unlock 'l'|
\end{Verbatim}
 Hash locks should only be used for very short operations, since there is a
chance that two concurrently locked objects map to the same hash value,
leading to unnecessary contention. 

 Hash locks are unrelated to the locks used by the \texttt{atomic} statements and the \texttt{LOCK} (\ref{LOCK}) and \texttt{UNLOCK} (\ref{UNLOCK}) primitives. 

\subsection{\textcolor{Chapter }{HASH{\textunderscore}LOCK}}
\logpage{[ 11, 2, 1 ]}\nobreak
\hyperdef{L}{X8145561486573102}{}
{\noindent\textcolor{FuncColor}{$\triangleright$\enspace\texttt{HASH{\textunderscore}LOCK({\mdseries\slshape obj})\index{HASH{\textunderscore}LOCK@\texttt{HASH{\textunderscore}LOCK}}
\label{HASHuScoreLOCK}
}\hfill{\scriptsize (function)}}\\


 \texttt{HASH{\textunderscore}LOCK} obtains the read\texttt{\symbol{45}}write lock for the hash value associated
with \texttt{obj}. }

 

\subsection{\textcolor{Chapter }{HASH{\textunderscore}UNLOCK}}
\logpage{[ 11, 2, 2 ]}\nobreak
\hyperdef{L}{X7AD0136085D9CFE0}{}
{\noindent\textcolor{FuncColor}{$\triangleright$\enspace\texttt{HASH{\textunderscore}UNLOCK({\mdseries\slshape obj})\index{HASH{\textunderscore}UNLOCK@\texttt{HASH{\textunderscore}UNLOCK}}
\label{HASHuScoreUNLOCK}
}\hfill{\scriptsize (function)}}\\


 \texttt{HASH{\textunderscore}UNLOCK} releases the read\texttt{\symbol{45}}write lock for the hash value associated
with \texttt{obj}. }

 

\subsection{\textcolor{Chapter }{HASH{\textunderscore}LOCK{\textunderscore}SHARED}}
\logpage{[ 11, 2, 3 ]}\nobreak
\hyperdef{L}{X7BE2F84778C68D04}{}
{\noindent\textcolor{FuncColor}{$\triangleright$\enspace\texttt{HASH{\textunderscore}LOCK{\textunderscore}SHARED({\mdseries\slshape obj})\index{HASH{\textunderscore}LOCK{\textunderscore}SHARED@\texttt{HAS}\-\texttt{H{\textunderscore}}\-\texttt{L}\-\texttt{O}\-\texttt{C}\-\texttt{K{\textunderscore}}\-\texttt{S}\-\texttt{H}\-\texttt{ARED}}
\label{HASHuScoreLOCKuScoreSHARED}
}\hfill{\scriptsize (function)}}\\


 \texttt{HASH{\textunderscore}LOCK{\textunderscore}SHARED} obtains the read\texttt{\symbol{45}}only lock for the hash value associated
with \texttt{obj}. }

 

\subsection{\textcolor{Chapter }{HASH{\textunderscore}UNLOCK{\textunderscore}SHARED}}
\logpage{[ 11, 2, 4 ]}\nobreak
\hyperdef{L}{X877AEBBC8249DA2C}{}
{\noindent\textcolor{FuncColor}{$\triangleright$\enspace\texttt{HASH{\textunderscore}UNLOCK{\textunderscore}SHARED({\mdseries\slshape obj})\index{HASH{\textunderscore}UNLOCK{\textunderscore}SHARED@\texttt{HAS}\-\texttt{H{\textunderscore}}\-\texttt{U}\-\texttt{N}\-\texttt{L}\-\texttt{O}\-\texttt{C}\-\texttt{K{\textunderscore}}\-\texttt{S}\-\texttt{H}\-\texttt{ARED}}
\label{HASHuScoreUNLOCKuScoreSHARED}
}\hfill{\scriptsize (function)}}\\


 \texttt{HASH{\textunderscore}UNLOCK{\textunderscore}SHARED} releases the read\texttt{\symbol{45}}only lock for the hash value associated
with \texttt{obj}. }

 }

 
\section{\textcolor{Chapter }{Migration to the public region}}\label{Migration to the public region}
\logpage{[ 11, 3, 0 ]}
\hyperdef{L}{X834BA1408083326A}{}
{
  HPC\texttt{\symbol{45}}GAP allows migration of arbitrary objects to the public
region. This functionality is potentially dangerous; for example, if two
threads try resize a plain list simultaneously, this can result in memory
corruption. 

 Accordingly, such data should never be accessed except through operations that
protect accesses through locks, memory barriers, or other mechanisms. 

\subsection{\textcolor{Chapter }{MAKE{\textunderscore}PUBLIC}}
\logpage{[ 11, 3, 1 ]}\nobreak
\hyperdef{L}{X8011BF108795B266}{}
{\noindent\textcolor{FuncColor}{$\triangleright$\enspace\texttt{MAKE{\textunderscore}PUBLIC({\mdseries\slshape obj})\index{MAKE{\textunderscore}PUBLIC@\texttt{MAKE{\textunderscore}PUBLIC}}
\label{MAKEuScorePUBLIC}
}\hfill{\scriptsize (function)}}\\


 \texttt{MAKE{\textunderscore}PUBLIC} makes \mbox{\texttt{\mdseries\slshape obj}} and all its subobjects members of the public region. }

 

\subsection{\textcolor{Chapter }{MAKE{\textunderscore}PUBLIC{\textunderscore}NORECURSE}}
\logpage{[ 11, 3, 2 ]}\nobreak
\hyperdef{L}{X81142B5B86D8FE29}{}
{\noindent\textcolor{FuncColor}{$\triangleright$\enspace\texttt{MAKE{\textunderscore}PUBLIC{\textunderscore}NORECURSE({\mdseries\slshape obj})\index{MAKE{\textunderscore}PUBLIC{\textunderscore}NORECURSE@\texttt{MAK}\-\texttt{E{\textunderscore}}\-\texttt{P}\-\texttt{U}\-\texttt{B}\-\texttt{L}\-\texttt{I}\-\texttt{C{\textunderscore}}\-\texttt{N}\-\texttt{O}\-\texttt{R}\-\texttt{E}\-\texttt{C}\-\texttt{URSE}}
\label{MAKEuScorePUBLICuScoreNORECURSE}
}\hfill{\scriptsize (function)}}\\


 \texttt{MAKE{\textunderscore}PUBLIC{\textunderscore}NORECURSE} makes \mbox{\texttt{\mdseries\slshape obj}}, but not any of its subobjects members of the public region. }

 }

 
\section{\textcolor{Chapter }{Memory barriers}}\label{Memory barriers}
\logpage{[ 11, 4, 0 ]}
\hyperdef{L}{X81C8B4FA86E9DFD9}{}
{
  The memory models of some processors do no guarantee that read and writes
reflect accesses to main memory in the same order in which the processor
performed them; for example, code may write variable \texttt{v1} first, and \texttt{v2} second; but the cache line containing \texttt{v2} is flushed to main memory first so that other processors see the change to \texttt{v2} before the change to \texttt{v1}. 

 Memory barriers can be used to prevent such
counter\texttt{\symbol{45}}intuitive reordering of memory accesses. 

\subsection{\textcolor{Chapter }{ORDERED{\textunderscore}WRITE}}
\logpage{[ 11, 4, 1 ]}\nobreak
\hyperdef{L}{X7C15E19D85AD81E2}{}
{\noindent\textcolor{FuncColor}{$\triangleright$\enspace\texttt{ORDERED{\textunderscore}WRITE({\mdseries\slshape expr})\index{ORDERED{\textunderscore}WRITE@\texttt{ORDERED{\textunderscore}WRITE}}
\label{ORDEREDuScoreWRITE}
}\hfill{\scriptsize (function)}}\\


 The \texttt{ORDERED{\textunderscore}WRITE} function guarantees that all writes that occur prior to its execution or
during the evaluation of \mbox{\texttt{\mdseries\slshape expr}} become visible to other processors before any of the code executed after. 

 Example: 
\begin{Verbatim}[commandchars=!@|,fontsize=\small,frame=single,label=Example]
  !gapprompt@gap>| !gapinput@y:=0;; f := function() y := 1; return 2; end;;|
  !gapprompt@gap>| !gapinput@x := ORDERED_WRITE(f());|
  2
\end{Verbatim}
 Here, the write barrier ensure that the assignment to \texttt{y} that occurs during the call of \texttt{f()} becomes visible to other processors before or at the same time as the
assignment to \texttt{x}. 

 This can also be done differently, with the same semantics: 
\begin{Verbatim}[commandchars=!@|,fontsize=\small,frame=single,label=Example]
  !gapprompt@gap>| !gapinput@t := f();; # temporary variable|
  !gapprompt@gap>| !gapinput@ORDERED_WRITE(0);; # dummy argument|
  !gapprompt@gap>| !gapinput@x := t;|
  2
\end{Verbatim}
 }

 

\subsection{\textcolor{Chapter }{ORDERED{\textunderscore}READ}}
\logpage{[ 11, 4, 2 ]}\nobreak
\hyperdef{L}{X7C4DDAEF7A900359}{}
{\noindent\textcolor{FuncColor}{$\triangleright$\enspace\texttt{ORDERED{\textunderscore}READ({\mdseries\slshape expr})\index{ORDERED{\textunderscore}READ@\texttt{ORDERED{\textunderscore}READ}}
\label{ORDEREDuScoreREAD}
}\hfill{\scriptsize (function)}}\\


 Conversely, the \texttt{ORDERED{\textunderscore}READ} function ensures that reads that occur before its call or during the
evaluation of \mbox{\texttt{\mdseries\slshape expr}} are not reordered with respects to memory reads occurring after it. }

 }

 
\section{\textcolor{Chapter }{Object manipulation}}\label{Object manipulation}
\logpage{[ 11, 5, 0 ]}
\hyperdef{L}{X8217AB0984FDE424}{}
{
  There are two new functions to exchange a pair of objects. 

\subsection{\textcolor{Chapter }{SWITCH{\textunderscore}OBJ}}
\logpage{[ 11, 5, 1 ]}\nobreak
\hyperdef{L}{X84644A1284D0AB35}{}
{\noindent\textcolor{FuncColor}{$\triangleright$\enspace\texttt{SWITCH{\textunderscore}OBJ({\mdseries\slshape obj1, obj2})\index{SWITCH{\textunderscore}OBJ@\texttt{SWITCH{\textunderscore}OBJ}}
\label{SWITCHuScoreOBJ}
}\hfill{\scriptsize (function)}}\\


 \texttt{SWITCH{\textunderscore}OBJ} exchanges its two arguments. All variables currently referencing \mbox{\texttt{\mdseries\slshape obj1}} will reference \mbox{\texttt{\mdseries\slshape obj2}} instead after the operation completes, and vice versa. Both objects stay
within their previous regions. 
\begin{Verbatim}[commandchars=!@|,fontsize=\small,frame=single,label=Example]
  !gapprompt@gap>| !gapinput@a := [ 1, 2, 3];;|
  !gapprompt@gap>| !gapinput@b := [ 4, 5, 6];;|
  !gapprompt@gap>| !gapinput@SWITCH_OBJ(a, b);|
  !gapprompt@gap>| !gapinput@a;|
  [ 4, 5, 6 ]
  !gapprompt@gap>| !gapinput@b;|
  [ 1, 2, 3 ]
\end{Verbatim}
 The function requires exclusive access to both objects, which may necessitate
using an atomic statement, e.g.: 
\begin{Verbatim}[commandchars=!@|,fontsize=\small,frame=single,label=Example]
  !gapprompt@gap>| !gapinput@a := ShareObj([ 1, 2, 3]);;|
  !gapprompt@gap>| !gapinput@b := ShareObj([ 4, 5, 6]);;|
  !gapprompt@gap>| !gapinput@atomic a, b do SWITCH_OBJ(a, b); od;|
  !gapprompt@gap>| !gapinput@atomic readonly a do Display(a); od;|
  [ 4, 5, 6 ]
  !gapprompt@gap>| !gapinput@atomic readonly b do Display(b); od;|
  [ 1, 2, 3 ]
\end{Verbatim}
 }

 

\subsection{\textcolor{Chapter }{FORCE{\textunderscore}SWITCH{\textunderscore}OBJ}}
\logpage{[ 11, 5, 2 ]}\nobreak
\hyperdef{L}{X7AA8AD9E7B28961C}{}
{\noindent\textcolor{FuncColor}{$\triangleright$\enspace\texttt{FORCE{\textunderscore}SWITCH{\textunderscore}OBJ({\mdseries\slshape obj1, obj2})\index{FORCE{\textunderscore}SWITCH{\textunderscore}OBJ@\texttt{FOR}\-\texttt{C}\-\texttt{E{\textunderscore}}\-\texttt{S}\-\texttt{W}\-\texttt{I}\-\texttt{T}\-\texttt{C}\-\texttt{H{\textunderscore}OBJ}}
\label{FORCEuScoreSWITCHuScoreOBJ}
}\hfill{\scriptsize (function)}}\\


 \texttt{FORCE{\textunderscore}SWITCH{\textunderscore}OBJ} works like \texttt{SWITCH{\textunderscore}OBJ} (\ref{SWITCHuScoreOBJ}), except that it can also exchange objects in the public region: 
\begin{Verbatim}[commandchars=!@|,fontsize=\small,frame=single,label=Example]
  !gapprompt@gap>| !gapinput@a := ShareObj([ 1, 2, 3]);;|
  !gapprompt@gap>| !gapinput@b := MakeImmutable([ 4, 5, 6]);;|
  !gapprompt@gap>| !gapinput@atomic a do FORCE_SWITCH_OBJ(a, b); od;|
  !gapprompt@gap>| !gapinput@a;|
  [ 4, 5, 6 ]
\end{Verbatim}
 This function should be used with extreme caution and only with public objects
for which only the current thread has a reference. Otherwise, undefined
behavior and crashes can result from other threads accessing the public object
concurrently. }

 }

 }

 \def\indexname{Index\logpage{[ "Ind", 0, 0 ]}
\hyperdef{L}{X83A0356F839C696F}{}
}

\cleardoublepage
\phantomsection
\addcontentsline{toc}{chapter}{Index}


\printindex

\newpage
\immediate\write\pagenrlog{["End"], \arabic{page}];}
\immediate\closeout\pagenrlog
\end{document}
